\documentclass[11pt]{article}

\usepackage[utf8]{inputenc}
\usepackage[T1]{fontenc}
\usepackage{lmodern}
\usepackage[margin=1in]{geometry}
\usepackage{amsmath,amssymb}
\usepackage{booktabs}
\usepackage{array}
\usepackage{longtable}
\usepackage{tabularx}
\usepackage{graphicx}
\usepackage{caption}
\usepackage{enumitem}
\usepackage{hyperref}
\usepackage{url}
\newcolumntype{P}[1]{>{\raggedright\arraybackslash}p{#1}}
\captionsetup{font=small,labelfont=bf}
\setlength{\emergencystretch}{2em}
\hypersetup{
  colorlinks=true,
  linkcolor=blue,
  citecolor=blue,
  urlcolor=blue,
  pdftitle={Cosmo-Local Credit (CLC) White Paper},
  pdfauthor={William O. Ruddick and Mohamed Sohail}
}

\title{Cosmo-Local Credit (CLC): A Network for Routing Credit, Settling Commitments, and Financing a Healthy Cosmo-Local Economy}
\author{William O. Ruddick \\ Mohamed Sohail \\ Grassroots Economics Foundation \\ \texttt{info@grassecon.org}}
\date{Version 0.7 --- 24 September 2026}

\begin{document}
\maketitle

\begin{abstract}
Cosmo-Local Credit (CLC) is a pattern for network clearing houses and governance layers built around Commitment Pooling Protocol (CPP) curation markets. Commitment Pools (CPs) coordinate redeemable commitments (vouchers) using four pooling functions: \textbf{Curation} (Commitment Registry), \textbf{Valuation} (Value Index Registry), \textbf{Limitation} (Swap Limiter), and \textbf{Exchange} (Vault/Fee Registry: fee policy + custody/settlement). A CLC deployment can publish a registry of Commitment Pools and may add routing, monitoring, liquidity support, service fees, reserves, guarantees, or insurance under expressly published policies.

Cosmo-local means we share open standards and software globally, while keeping issuance, redemption, governance, and real-world accountability local.

CPP curation markets: independent CPs list (``curate'') which vouchers they accept and publish their own value indices, limits, fees, inventory rules, and any applicable guarantee structures. A CLC-compatible registry or service layer can provide network routing, shared standards, and optional insurance layers, but does not automatically guarantee any specific pool or voucher unless an identified party expressly assumes that obligation in published terms.
\end{abstract}

\noindent\textbf{Status:} Draft for review and discussion only. This document is not investment advice and does not constitute an offer to sell or a solicitation to buy any security or financial instrument.

\tableofcontents
\newpage

\section*{Front Matter and Key Concepts}
\addcontentsline{toc}{section}{Front Matter and Key Concepts}

\textbf{Version:} 0.7

\textbf{Date:} 24 September 2026

\textbf{Distribution:} Limited-circulation draft for review and discussion only. This document is not investment advice and does not constitute an offer to sell or a solicitation to buy any security or financial instrument. Subject to change.

\textbf{Framing note:} This paper combines deployed capabilities with proposed governance and service models. The historic Sarafu Network service transitioned to Cosmo-Local Credit, but that transition does not assert that every account, wallet, token, pool, balance, report, or transaction was migrated. Other registries may be managed by different institutions or governance bodies.

\textbf{Intended audience:} Pool builders/stewards, protocol engineers/auditors, impact-liquidity providers and mandate funders, governance designers, and legal/compliance reviewers.

\section*{Current Deployment}
\addcontentsline{toc}{section}{Current Deployment}

As of 24 September 2026, the public progressive web app at \texttt{cosmolocal.credit} provides account and Wallet access, a public market, and interfaces for creating or using Tokens, Vouchers, Commitment Pools, and Swaps on Gnosis Chain. Documentation is published at \texttt{docs.cosmolocal.credit}. Available actions remain subject to the App interface, the relevant smart contracts, network state, inventory, limits, fees, and the terms published by each issuer and Pool Steward.

Grassroots Economics Foundation (GEF) operates the public App and supporting infrastructure. Operating the interface does not by itself make GEF an issuer, Pool Steward, custodian, lender, borrower, broker, guarantor, redeemer, adviser, insurer, or party to obligations between users. GEF does not guarantee the legality, value, liquidity, convertibility, redemption, or performance of any Token, Voucher, Pool, or Swap. Users remain responsible for reviewing transaction terms and complying with the laws that apply to them. Use of the App is governed by the \href{https://docs.cosmolocal.credit/governance/terms}{Terms of Service}.

The historic Sarafu Network service preceded this deployment and transitioned to Cosmo-Local Credit. Historical metrics in this paper describe that earlier service and are not current App metrics or evidence that every historical account, asset, record, or obligation migrated.

\subsection*{Pool Sovereignty \& Forkability.}
\addcontentsline{toc}{subsection}{Pool Sovereignty \& Forkability.}

As a design principle, each Commitment Pool's Steward controls its listings, limits, fees, any guarantees, and routing preferences within the deployed contracts and applicable terms. A compatible Pool should be able to leave a discovery registry without changing valid user balances or commitments. ``Canonical'' registries and routers are intended as convenience layers for discovery and shared standards, not as a monopoly.

The broader CLC technical design is intentionally \textbf{forkable}. Compatible deployments should use reproducible open-source software, exportable registries, and a documented ``fork \& migrate'' procedure (see \S{}11.5) so communities and operators can credibly exit if governance is captured or drifts from shared values.

\subsection*{Key Concepts:}
\addcontentsline{toc}{subsection}{Key Concepts:}

\begin{itemize}[leftmargin=*]
\item \textbf{Confederation note:} CLC is designed so that many independent CPP networks can coexist (each with its own registry roots, compliance policies, and insurance scope) while still routing to one another through multi-profile discovery. Canonical registries/routers are a convenience layer, not a monopoly; credible exit is preserved by fork-and-repoint (see Sections 5.2a, 8.1, and 11.5).
\begin{itemize}[leftmargin=*]
\item A ``profile'' is a named bundle of registry roots + routing policies (and optional compliance/insurance scope) that a community can choose as its discovery and safety baseline.
\item Registry management and service-fee recipients are deployment-specific. Any institution or governance body can set up a Commitment Pool registry and charge transparent fees for services such as routing, clearing, monitoring, liquidity support, or insurance coordination. A network may be stewarded by a nonprofit foundation, cooperative, holding company, public agency, community group, multisig, service operator, or another accountable structure. GEF operates the current public Cosmo-Local Credit App; the historic Sarafu Network was an earlier GEF-operated service.
\end{itemize}
\item \textbf{Cosmo-Local: Cosmopolitan-localism} links local communities through shared global infrastructures while keeping production, redemption, and governance local. In practice: we share ``light'' resources globally (standards, software, knowledge, registries), while keeping ``heavy'' realities local (relationships, fulfillment, ecology, material production, and legal accountability). This aligns with ``Small, Local, Open and Connected'' (SLOC) and ``design global, manufacture local'' (DGML): share what's light; keep what's heavy local.  See \href{https://en.wikipedia.org/wiki/Cosmopolitan_localism}{Cosmopolitan\_localism}
\item \textbf{Cosmo-Local Credit (CLC):} CLC applies this to obligations: commitments are issued and redeemed locally, while any guarantee must be expressly assumed and disclosed. Compatible commitments may be discovered and routed across trusted Pools using common registries, transaction records, and configured safety constraints.
\begin{itemize}[leftmargin=*]
\item Example: a clinic voucher remains a local promise, but it can be swapped into food or transport credits by routing through trusted pools, with receipts showing what happened and why.
\end{itemize}
\item \textbf{Cosmo-Local Economy:} A cosmo-local economy is a federation of local markets that interoperate and clear obligations without a central monopoly. Communities can choose their routers/registries, publish their own policies, and retain credible exit (forkability) if governance is captured.
\item The digital layer is a \textbf{shared memory + coordination tool}, not a replacement for in-person relationships or local accountability. We explicitly design against techno-solutionism and capture risks by prioritizing: credible exit/forkability, local sovereignty, transparent receipts, timelocks/quorums, bounded limits/reserves, and low-tech access (assisted flows).. Where proof matters, vouchers may require simple evidence (signatures/photos/device proofs) governed locally.
\item \textbf{Producer Credit (optional, expressly documented):} A separately established credit facility may accept curated vouchers as collateral or a repayment instrument. It is a loan only when the transaction is expressly presented as a repayable advance and states the parties, principal or advance, repayment amount, fees or interest, due date, collateral treatment, assignment, default, remedies, and discharge evidence. An ordinary Swap, the direction of assets in a Swap, or redemption of a Voucher does not by itself create, transfer, repay, or discharge a loan.
\end{itemize}
\subsection*{Quick Facts}
\addcontentsline{toc}{subsection}{Quick Facts}

The following points summarize the full CLC design. Governance tokens, multi-hop routing, shared insurance, netting, liquidity mandates, and fee-credit mechanisms are proposed or deployment-specific unless expressly identified as active.

\begin{enumerate}[leftmargin=*]
\item \textbf{What this is:} CLC describes clearing-network patterns for \textbf{redeemable commitments} (think: tokens, \textbf{digital gift-cards / service credits / delivery claims}) that can be exchanged across curated markets.
\begin{enumerate}[leftmargin=*]
\item Not 1-to-1 barter: barter is usually a direct trade between two people. A Commitment Pool enables multi-party exchange - you can contribute value to the pool and later redeem from the pool's shared inventory. It's multilateral barter through a shared intermediary (the pool) - CLC expands that by routing between a network of pools..
\end{enumerate}
\item \textbf{Values:}
\begin{enumerate}[leftmargin=*]
\item \textbf{Care for People}-- collaboration and vision driven care for oneself and others' well-being and happiness
\item \textbf{Care for the Environment}- support environmental protection and regeneration, minimize use of finite resources for economic activity, ecosystem management approach to farming and business development.
\item \textbf{Fairness}- fair and secure access to Instruments, land, resources, knowledge \& care for members from different backgrounds, age, gender and religion.
\item \textbf{Reciprocity}-- mutual sharing of risk, cost and surplus
\item \textbf{Non-Dominance}- no person or association to have dominant rights over another person or association's resources eg. data, finances, intellect, materials and freedom.
\item \textbf{Resilience}- capacity to prepare for, address and adapt to economic, political, climate and other events in order to ensure sustainable community based systems/commons.
\end{enumerate}
\item \textbf{How we reach sustainability / How fees are generated:} Fees come from \textbf{settlement activity} (people swapping, redeeming, and fulfilling goods/services), not price speculation. Where enabled, service fees on routed swaps/settlements across participating CPs fund shared services. These fees may be a percentage of pool fees, a routing/clearing fee, or another published fee schedule.
\begin{enumerate}[leftmargin=*]
\item Opt-in, not extraction: service fees apply only to pools and routes that opt into a registry, routing profile, or shared service layer. Pools remain sovereign and can exit canonical registries/routers at any time while continuing to operate locally.
\item Not trickle-down: pooled fees are allocated by a published Waterfall - insurance targets, core operations, and liquidity/off-ramp mandates come first. Only after those priorities are met can a capped ``fee-access budget'' be published.
\end{enumerate}
\item \textbf{How Liquidity Providers could participate in the proposed model:}
\begin{enumerate}[leftmargin=*]
\item Provide liquidity by seeding designated Commitment Pools and, only under a separately adopted program, become eligible for policy-gated access to pooled fees or CLC governance tokens via the Impact Seeding program (\S{}7.2.3).
\item Support liquidity mandates (policy-directed endowments; reporting-heavy; impact-first).
\item Under that proposed model, participants could lock or escrow CLC to receive stCLC voting power.
\item An implemented epoch policy could also mint sCLC as a bounded authorization or incentive. Where expressly enabled, sCLC could provide capped Swap access into designated fee-holding Pools after the Waterfall. It would provide no dividends, profit share, or residual rights.
\end{enumerate}
\item \textbf{What a future LP program could fund:} Liquidity endowments (e.g., stable cash-equivalents) intended to help real-world Vouchers be exchanged and redeemed, subject to the program's published risks and terms.
\item \textbf{Downside controls:} Limits and inventory checks can reduce certain risks. Reserves, guarantees, insurance, timelocks, dashboards, and credible exit are optional controls and protect users only when implemented, funded, and expressly disclosed; they do not eliminate loss.
\item \textbf{What a deployment may publish:} Dashboards and receipts for volume, settlement rate, redemption performance, reserves, incidents, fee flows, and governance changes, plus the KPIs used for any mandates or insurance targets.
\item What moves? Redeemable commitments (vouchers).
\item What's the goal? Increase velocity of settlement of real-world obligations - while maintaining the CLC values.
\begin{enumerate}[leftmargin=*]
\item We treat the digital layer as shared memory and coordination (not as a substitute for human relationships) and design explicitly against capture (timelocks, transparency, limits, and credible exit).
\end{enumerate}
\item \textbf{Who is responsible for what?} Issuers are responsible for honoring their published Voucher commitments. Pool Stewards are responsible for the Pool rules and any guarantees they expressly publish. A CLC deployment may provide routing, standards, governance, or an optional insurance policy, but neither CLC nor GEF automatically guarantees any Pool, Voucher, value, liquidity, or redemption. Any guarantee or coverage must identify the responsible party, scope, funding, caps, exclusions, claim process, and applicable terms.
\item Proposed participant access? If separately implemented and enabled by policy, eligible staked-CLC participants could receive \textbf{time-bounded Swap access} through sCLC to a defined portion of pooled fees under disclosed caps and windows; governance could set this access to \textbf{zero} in any epoch.
\item Safety? A CP can configure per-Voucher or rolling-window limits and may publish reserve, guarantor, insurance, or circuit-breaker policies. Users must verify which controls actually apply to a particular Pool and transaction.
\end{enumerate}
Reader Map (where to look)

\begin{itemize}[leftmargin=*]
\item If you care about confederation \& interoperability: see \S{}5.2a, \S{}8.1, and \S{}11.5.
\item If you care about guarantees / insurance boundaries: see \S{}11.3.
\item If you care about LP economics and fee flow: see \S{}7.4, \S{}9, and \S{}12.
\item If you care about risk controls: see \S{}6 and \S{}10.
\item If you care about forkability and credible exit: see \S{}8.2 and \S{}11.5.
\end{itemize}

\section*{Executive Summary}
\addcontentsline{toc}{section}{Executive Summary}

Modern financial systems excel at trading volatile assets but struggle to finance real production, community resilience, and long-term commitments. The \textbf{Commitment Pooling Protocol (CPP)} offers an alternative: a simple, extensible protocol for issuing, routing, and settling redeemable commitments (claims on future goods, services, and labor) within and across communities. Issuers are responsible for their published Voucher commitments; Pools may add expressly disclosed guarantees; neither CLC nor GEF is a universal guarantor.

The historic \textbf{Sarafu Network} service used CPP concepts to coordinate community Vouchers, savings groups, mutual-aid systems, and production commitments. That service later transitioned to the Cosmo-Local Credit PWA. This was a service transition, not a representation that every historical account, Wallet, Token, Pool, balance, report, transaction, or obligation migrated.

\subsection*{Historical Sarafu Network snapshot (Celo activity from 5 July 2023 through 20 July 2025):}
\addcontentsline{toc}{subsection}{Historical Sarafu Network snapshot (Celo activity from 5 July 2023 through 20 July 2025):}

\textbf{Historical on-chain source:} \href{https://dune.com/grassrootseconomics/sarafu-network}{Dune Analytics dashboard}

\begin{itemize}[leftmargin=*]
\item 26,367 users
\item 285,197 peer-to-peer exchanges
\item 188 unique active commitment pools
\item 745 unique active vouchers
\item \$320,692 pool swap volume
\item 899 impact reports published (\href{https://sarafu.network/reports}{historical reports archive})
\end{itemize}
These figures preserve the provenance of the earlier service. They are not current Cosmo-Local Credit usage figures and do not establish continuity of any particular user or asset.

The historical experience motivated research into optional liquidity and governance layers that could:

\begin{enumerate}[leftmargin=*]
\item Inject liquidity across pools,
\item Coordinate accountable governance as the network scales,
\item Coordinate expressly scoped and funded shared insurance policies and incident runbooks,
\item Provide bounded, policy-gated fee-credit (ex-post swap access to a defined portion of pooled fees) for those who underwrite risk and coordination.
\end{enumerate}
The broader \textbf{Cosmo-Local Credit (CLC)} design explores how to fulfill these roles. These proposed layers are distinct from the capabilities currently exposed by the public PWA.

CLC is designed as a win-win routing layer across curated commitment markets:

\begin{enumerate}[leftmargin=*]
\item Pool Stewards curate Voucher listings and publish rules and any guarantees they choose to assume; their performance can become discoverable and comparable.
\item Where a separate credit facility is expressly documented, lenders and liquidity providers may finance real production while accepting redeemable Vouchers as collateral or a repayment instrument.
\item Under those transaction-specific credit terms, producers or borrowers may repay through agreed cash or in-kind performance.
\item Consumers can browse curated markets and redeem Vouchers under issuer terms. An ordinary Swap, its asset direction, or a Voucher redemption does not by itself create or discharge a loan.
\item Governance participants direct liquidity mandates to increase settlement velocity and may receive bounded, policy-gated access to eligible fee assets only as defined by published policy.
\end{enumerate}
The proposed CLC design includes a \textbf{network governance and liquidity coordination layer (including a CLC token design)} intended to align liquidity providers, Pool creators, communities, and Stewards around one shared goal:

\textbf{Increase the velocity of settlement of real-world commitments while preserving care, fairness, and resilience.}

\textbf{A worst-case governance scenario:} A hostile actor accumulates CLC voting power (e.g., via public markets) and attempts to redirect fee flows, weaken curation standards, or force liquidity mandates that harm communities.

This design therefore treats anti-capture and credible exit (forkability) as first-class safety properties: (i) time-locked and multi-threshold governance for critical parameters, (ii) voting power that requires lockups (no instant governance via spot purchases), (iii) transparent delegation and conflict-of-interest rules, and (iv) a documented fork-and-migrate process so pools and communities can exit if governance is captured (see \S{}11.5).



\section*{1. Commitment Pooling Protocol (CPP): The Core Primitive}
\addcontentsline{toc}{section}{1. Commitment Pooling Protocol (CPP): The Core Primitive}

\textbf{Mental model:} A Commitment Pool is like a small, governed \textbf{clearing house} for community ``gift cards'' (Vouchers). People deposit Vouchers or other accepted assets, exchange them under published rules, and redeem Vouchers for the goods or services described by their issuers. Configured contracts can enforce limits and record transactions. Any guarantee, remedy, or cross-Pool route exists only where it is expressly implemented and disclosed.

CPP is a protocol for coordinating value using \textbf{commitments}. Commitments are the economy; CLC makes them open and routable. CPP is described in the book, \href{https://willruddick.substack.com/p/grassroots-economics-the-book-is}{Grassroots Economics: Reflection and Practice.}

\subsection*{1.1 What is a Commitment?}
\addcontentsline{toc}{subsection}{1.1 What is a Commitment?}

A commitment is a clearly defined promise of future delivery - e.g., maize, transport services, labor hours, storage, or currency redemption. These commitments are represented as \textbf{vouchers}, which function economically like \textbf{pre-paid delivery claims} (similar to gift cards / service credits). (further defined in the section of voucher schemas).

\subsection*{1.2 What is a Commitment Pool?}
\addcontentsline{toc}{subsection}{1.2 What is a Commitment Pool?}

Roles:

\begin{itemize}[leftmargin=*]
\item Pool steward: decides listings, values, fees, limits, guarantees, and pauses for a specific pool.
\item CLC shared services: standards/registries, routing policies, monitoring, liquidity/off-ramps, and optional insurance.
\item Router/operator: finds and executes paths across pools; can prefer safer profiles and degrade toxic routes.
\item Guarantor (optional): backs specific vouchers/pools with an explicit remedy if commitments fail.
\end{itemize}
A Commitment Pool is a stewarded contract suite that implements four interfaces:

\begin{itemize}[leftmargin=*]
\item \textbf{Curation}: Registers acceptable vouchers (Commitment (token) Registry),
\item \textbf{Valuation}: Maintains a value index (Value Index Registry),
\begin{itemize}[leftmargin=*]
\item The Value Index is the Pool's published valuation reference for quoting one accepted asset relative to another. A reference is not a guarantee of market price, redemption value, liquidity, or convertibility. Pools using different valuation methods may interoperate where each Pool and route can quote and enforce compatible rules.
\end{itemize}
\item \textbf{Limitation}: Enforces swap limits and, only for an expressly designated credit facility, any configured credit limits (Swap Limiter),
\item \textbf{Exchange (Vault/Fee Registry):} Configures fees and custodies assets; executes seed/swap only if listed, valued by the pool index, within limits, and in stock; emits receipts for every action.
\end{itemize}
Each Pool can behave like a \textbf{locally governed clearing house}, stewarded by an individual, cooperative, community group, public agency, federation, multisig, service operator, or another accountable structure.

\textbf{Big idea:} We are already doing commitment pooling all the time: wages, rent, invoices, loans, warranties, memberships, and mutual aid are all promises that get trusted, netted, and settled. Today that pooling is mostly closed, opaque, and permissioned (inside institutions and platforms) so commitments can't easily connect or route beyond their enclosures. \textbf{CLC proposes to make the underlying protocol open and interoperable}, so commitments can be published, pooled, and routed across communities and markets - by anyone.

\textbf{Why this matters:} pools can be compared and risk-rated because their listings, limits, fees, reserves, and guarantees are explicitly published.

Minimal Swap Logic (canonical)

This is reference logic for a compatible Pool. A particular deployment or user interface may expose only a subset, such as direct Swaps; multi-hop routing and alternative settlement mechanisms are optional capabilities rather than guaranteed paths.

\begin{enumerate}[leftmargin=*]
\item Listed?: Input and output vouchers must be listed in the Commitment Registry.
\item Price?: Compute input$\to$output via the Value Index (with fee preview).
\item Limits?: Enforce Swap Limiter windows/caps for each voucher (and account/global if configured).
\item Exchange:
\begin{enumerate}[leftmargin=*]
\item (fees): Apply Vault/Fee Registry fee rules (pair-specific OK); preview and emit on receipt.
\item (inventory): Verify Vault inventory for the outgoing voucher.
\end{enumerate}
\item Transfer \& Log: Move in/out, update registries, and emit an immutable receipt that links the quote to the receipt.
\item Multi-hop?: Repeat hop-by-hop; atomic if supported, else HTLC/escrow with explicit abort paths.
\end{enumerate}
\textbf{In other words:} a pool only lets people exchange vouchers that are (a) approved, (b) priced by the pool's published index, (c) within safety caps, and (d) actually in stock - then it issues a receipt.

\subsection*{1.3 Why CPP is Different from Traditional Crypto Decentralized Exchanges (DEXs)}
\addcontentsline{toc}{subsection}{1.3 Why CPP is Different from Traditional Crypto Decentralized Exchanges (DEXs)}

Most decentralized exchanges (automated trading pools) rely on bilateral token pairs, continuous curves, and volatility-driven fees. CPP supports those use cases \emph{and} enables low-frequency, high-impact coordination:

\begin{itemize}[leftmargin=*]
\item Community savings (VSLAs),
\item Production financing,
\item Mutual credit and mutual aid,
\item Insurance and guarantees,
\item Lending against real output,
\item Settlement support for separately documented personal or institutional debts.
\item Portfolio-directed liquidity (e.g. curated pools for ecosystem services, humanitarian support, and health \& wellness)
\end{itemize}
\emph{CPP is optimized for \textbf{fulfillment} and \textbf{auditable receipts}, not speculative churn.}



\section*{2. The Accounting Shift: From Assets to Trust}
\addcontentsline{toc}{section}{2. The Accounting Shift: From Assets to Trust}

Traditional finance begins with \textbf{Assets - Liabilities = Equity}. CPP reframes this for a commitment economy:

\begin{itemize}[leftmargin=*]
\item \textbf{Acceptance Capacity}: How much of an issuer's commitments a Pool or network is still willing to accept.
\item \textbf{Outstanding Commitments}: Issued promises that remain unfulfilled and are held by others.
\item \textbf{Backing Capacity}: Your real ability to honor those promises.
\end{itemize}
\subsection*{Commitment--Capacity Identity:}
\addcontentsline{toc}{subsection}{Commitment--Capacity Identity:}

Acceptance Capacity - Outstanding Commitments = Remaining Capacity

This conceptual identity can inform Pool risk limits: unlimited issuance does not imply unlimited acceptance. It is not a balance-sheet identity or a statement that every Voucher is legally a debt or loan.

\textbf{Example:} A transporter issues ``100 rides'' in Vouchers. A Pool accepts at most 40 rides worth at a time. If 15 accepted rides remain outstanding, the Pool has capacity to accept up to 25 additional rides under that policy. The calculation does not itself create a loan or guarantee fulfillment.



\section*{3. Velocity of Settlement: Why Liquidity Providers Should Care}
\addcontentsline{toc}{section}{3. Velocity of Settlement: Why Liquidity Providers Should Care}

We reframe ``velocity of money'' as velocity of settlement: how quickly outstanding promises move from owed to fulfilled.

For a given voucher type j across a network of CPs, define:

\begin{enumerate}[leftmargin=*]
\item D\_j = total outstanding Voucher commitments valued in a common index
\item S\_j = total value of settlements (redemptions routed through CPs) per period
\end{enumerate}
Then the network settlement velocity of voucher j is:

\[
V_j^{\mathrm{network}} = \frac{S_j}{D_j}
\]

This is a flow/stock ratio: how many units of fulfillment flow pass through the network per unit of outstanding Voucher commitments. We will expand this below to a federation of CPs.

\textbf{Key insight:} an ordinary Swap changes who holds assets; its direction does not create or repay a loan. Redemption can fulfill an issuer's Voucher commitment. It reduces a separate debt only when express, transaction-specific credit terms say that redemption constitutes repayment and provide evidence of discharge. Liquidity that increases routing capacity may increase fulfillment velocity and fee volume.

\textbf{Plain language:} If vouchers get redeemed quickly, more real trade flows through the network. More flow $\to$ more fee events $\to$ more sustainable fee pooling.

What This Is / Isn't

\begin{itemize}[leftmargin=*]
\item Not an AMM for speculative pairs. CPP values and limits are policyful and capacity-aware.
\item Not necessarily a bank deposit, payment instrument, security, or other regulated product; classification depends on the asset, terms, activities, and jurisdiction.
\item Does not make unbounded issuance safe. Pool limits and inventory checks can constrain eligible Swaps, but do not constrain an issuer's total supply unless expressly configured to do so.
\item Is a clearing-network design for redeemable commitments with transaction records and any recourse that the responsible parties expressly publish.
\item Can support a separately documented producer-credit facility in which Vouchers serve as collateral or a repayment instrument. No ordinary Swap or redemption is a loan or repayment merely because of asset direction.
\end{itemize}


\section*{4. Reusable Forward-Style Collateral}
\addcontentsline{toc}{section}{4. Reusable Forward-Style Collateral}

An expressly documented credit facility may accept \textbf{fungible, transferable Vouchers} (such as service credits or production forwards) as collateral or as a repayment instrument. If it does:

\begin{itemize}[leftmargin=*]
\item Collateral may become more discoverable across participating Pools,
\item Additional lawful redemption paths may improve fulfillment options,
\item Limits, inventory, value, liquidity, and default risks remain, and
\item No route, redemption, repayment, value, or recovery is guaranteed.
\end{itemize}
\subsection*{4.1 Producer Credit Loop (Loan Repayment via Curated Vouchers)}
\addcontentsline{toc}{subsection}{4.1 Producer Credit Loop (Loan Repayment via Curated Vouchers)}

CLC can support producer credit only when the parties establish a distinct facility and expressly present the transaction as a repayable advance. Its transaction terms must identify the creditor and debtor and state the advance and repayment amounts, due date, interest or fees, collateral use, changes in the creditor or Holder, assignment, partial settlement, redemption valuation, evidence of discharge, default, and lawful remedies.

\begin{enumerate}[leftmargin=*]
\item A producer (or service provider) issues a voucher: a redeemable claim on their future output (e.g., ``10 taxi rides,'' ``50kg maize,'' ``10 labor-hours'').
\item A Pool Steward may list that Voucher and publish the Pool's limits, fees, valuation, inventory rules, and any optional reserve or guarantee terms.
\item A lender or liquidity program separately provides an advance under written transaction terms and takes the producer's Vouchers as collateral or an agreed repayment instrument.
\item Holders may Swap or redeem those Vouchers under the applicable issuer and Pool terms.
\item A redemption fulfills the relevant Voucher commitment. It reduces the separate credit balance only to the extent, at the valuation, and upon the evidence specified in the credit terms.
\item Transaction records should distinguish Voucher fulfillment from loan accounting and identify any partial settlement, remaining amount, assignment, default, or discharge.
\end{enumerate}
This structure can permit agreed in-kind repayment while preserving separate records for the Voucher and credit obligations. It does not arise from an ordinary Swap, from which asset enters or leaves a Pool, or from redemption alone.

\textbf{Mini example:} A producer receives a \$1,000 advance under terms stating that verified redemption of specified maize Vouchers is credited against the repayment amount at a stated Full Value. A Holder's redemption fulfills the maize Voucher; the lender credits the amount only after receiving the evidence required by those terms. If the terms do not make that connection, buying, swapping, or redeeming the Voucher has no effect on any loan.


\section*{5. From Isolated Pools to a Federated Network}
\addcontentsline{toc}{section}{5. From Isolated Pools to a Federated Network}

CPs interoperate when they list the same vouchers. Routers move value across pools, respecting each hop's \textbf{value index, limits, fees, and inventory}. As pools proliferate:

\textbf{Status note:} Direct Pool interactions are available in the public PWA. Multi-hop routing, cross-profile routing, batch netting, shared Waterfall budgets, and sCLC mechanisms in this chapter are optional or proposed capabilities unless a deployment expressly identifies them as active.

\begin{itemize}[leftmargin=*]
\item Possible routing paths may multiply,
\item Fulfillment velocity may increase,
\item Fee volume may grow, and
\item Participants may gain access to inventories beyond a single Pool.
\end{itemize}
\textbf{Routing story:} A school accepts ``maize vouchers'' but needs ``transport vouchers.'' A router finds a path across pools that accept both. The swap clears only if each hop is within limits and inventory \ldots{} so the voucher reaches someone who can actually redeem it.

\subsection*{5.1 Velocity Multiplier}
\addcontentsline{toc}{subsection}{5.1 Velocity Multiplier}

When pools operate in isolation, each voucher can only settle within the small local circle that recognizes it. As pools are federated via a common protocol:

\begin{enumerate}[leftmargin=*]
\item More routes appear: a voucher can cross several pools to reach someone who can redeem it.
\item Credit is aggregated: multiple pools' acceptance capacity can support the same voucher type.
\item Netting surfaces: multi-lateral swaps reduce the need for bilateral matching.
\end{enumerate}
The result is higher settlement flow for a given stock of obligations: the same commitments can find fulfillment faster.

Formally, if we index pools by k:

\[
S_j = \sum_k S_{j,k}
\]

\[
D_j = \sum_k D_{j,k}
\]

Connectivity increases per-pool settlement flow without necessarily increasing per-pool obligation stock.

Thus settlement velocity for voucher j rises across the network as routing improves.

\subsection*{5.2 Routing as a Service: Pathfinding + Rebalancing (Liquidity-Saving)}
\addcontentsline{toc}{subsection}{5.2 Routing as a Service: Pathfinding + Rebalancing (Liquidity-Saving)}

In a CLC-compatible network, ``routing'' is not only a user-facing convenience (``swap voucher A for voucher B''). It is also a network liquidity service that increases settlement velocity by improving how inventory is distributed across pools and by surfacing multilateral netting opportunities.

Two routing modes:

\begin{enumerate}[leftmargin=*]
\item End-user routing (on-demand)
\end{enumerate}
Given (token\_in, token\_out, amount, constraints), the router discovers a multi-hop path across CPs. A route is valid only if each hop clears: (i) listing/registry checks, (ii) value index pricing, (iii) swap limiter windows/caps, (iv) fee rules, and (v) outgoing inventory availability. Execution is atomic where possible; otherwise HTLC/escrow is used with explicit abort paths.

\begin{enumerate}[leftmargin=*]
\item Pool rebalancing / batch netting (scheduled or threshold-triggered)
\end{enumerate}
Pools may opt-in to publish ``rebalance intents'' (or standing constraints): target inventory bands by voucher class, maximum deviation vs. the pool's value index, per-epoch caps, and allow/deny routing preferences. Routers/clearing agents then search for multilateral cycles and chains that:

(i) satisfy every hop's limits and inventories,

(ii) reduce inventory imbalance across pools, and

(iii) maximize total off-set value subject to policy constraints.

These cycles are executed as batch routes, producing receipts per hop (quote $\to$ receipt mapping).

\subsection*{Why this works}
\addcontentsline{toc}{subsection}{Why this works}

Obligation networks can contain cycles. When expressly compatible obligations are processed together, matching and fulfillment may require less external liquidity than sequential bilateral processing. This proposed ``cycle surfacing'' mechanism remains subject to every Pool's valuation, limits, fees, inventory, authorization, and applicable terms. It does not make an ordinary Swap a loan repayment.

\subsection*{Opt-in \& sovereignty note}
\addcontentsline{toc}{subsection}{Opt-in \& sovereignty note}

Rebalancing is never forced. A pool can be routable for end-users while disabling outbound rebalancing, or can enable only specific voucher classes, caps, and counterparties.

As routing and netting improve, the network can produce more fee events from real settlement activity; this can expand (policy-permitting) the fee-credit budget that defines sCLC swap-access capacity - without turning sCLC into an equity or profit instrument.

\subsection*{5.2a Confederation \& Interoperability}
\addcontentsline{toc}{subsection}{5.2a Confederation \& Interoperability}

CLC is designed as a confederation protocol: many independent networks can run their own registries, routers, and policy layers, while still routing to one another when they share compatible vouchers and standards. This is not a hub-and-spoke monopoly; it is a mesh of overlapping curations.

\subsection*{Interoperability incentive (why open source + forks help everyone):}
\addcontentsline{toc}{subsection}{Interoperability incentive (why open source + forks help everyone):}

\begin{itemize}[leftmargin=*]
\item Any fork/network that remains CPP-compatible can route to CLC pools (and CLC routers can route to theirs), increasing settlement paths, inventory reach, and real-world fulfillment velocity for all parties.
\item More interoperable networks $\to$ more routable paths $\to$ higher throughput and fee volume from real settlement activity (not speculation), benefiting LP programs and routing services across the confederation.
\item Open-source ``forkability'' reduces systemic risk: if any canonical registry/router becomes captured or degraded, communities can re-point or fork without bricking local economies.
\end{itemize}
\subsection*{Confederation mechanics (how it works):}
\addcontentsline{toc}{subsection}{Confederation mechanics (how it works):}

\begin{enumerate}[leftmargin=*]
\item Multi-profile discovery and explicit user/pool choice of registry roots (``network profiles'') (see \S{}8.1).
\item Reciprocal routing: confederated networks can publish mutual allowlists (registry roots / bridge adapters) with risk parameters (caps, escrow requirements, health-score thresholds).
\item Policy separation: each profile defines its own fee norms, insurance scope, and compliance hooks; routing across profiles must satisfy each hop's stated constraints and inventory.
\end{enumerate}
\subsection*{5.3. From Settlement Velocity to Fee Volume}
\addcontentsline{toc}{subsection}{5.3. From Settlement Velocity to Fee Volume}

Every routed swap or settlement can carry a small fee, analogous to an interchange fee in card networks.

Let:

\begin{enumerate}[leftmargin=*]
\item $\tau$ = average fee rate (e.g. 0.2\% per routed value unit)
\item D\_tot = ``how many promises exist'' (total outstanding redeemable obligations across vouchers, valued in the common index)
\item V(settlement) = aggregate settlement velocity across all vouchers
\end{enumerate}
Then approximate total fee revenue per period as:

\[
F \approx \tau \cdot V_{\mathrm{settle}} \cdot D_{\mathrm{tot}}
\]

D\_tot is ``how many promises exist'' (total outstanding redeemable obligations across vouchers), V\_settle is ``how fast they move,'' and $\tau$ is the effective service-fee rate per unit of routed value. Note that this may be a percentage of the fees that pool stewards charge.

\textbf{Rake-on-rake clarification (pool fees $\to$ network rake).} Pool stewards set a per-pool usage fee f\_p (as \% of value routed through that pool). CLC policy sets a rake share r\_p (as \% of that pool's collected fees). The effective network fee rate contributed by that pool is:

\[
\tau_p = f_p \cdot r_p
\]

The network-wide $\tau$ is the routed-value-weighted average across pools and routes (plus any router fees when applicable).

Convertibility constraint (cash-eligible vs. in-kind fees). Fees are collected in the same asset that moves through pools. If fees arrive as clinic credits or service vouchers, they can't directly pay for auditors, incident response, or insurance unless they are settled in-network or converted under policy.

Some fee assets are cash-equivalent/convertible (stables, major liquid tokens), while others are not (non-fiat-redeemable vouchers). Let $\chi$ be the share of total fee inflows that are cash-eligible/convertible after slippage and policy constraints. Define cash-usable fee revenue as:

\[
F_{\mathrm{cash}} \approx \chi \cdot F
\]

Eligibility \& conversion note: Fee inflows may include both cash-eligible assets (stables / major liquid tokens) and in-kind voucher assets. The protocol may convert allowlisted cash-eligible assets into stables/fiat when needed to meet insurance payouts, maintain off-ramp liquidity, and fund core operations - while prioritizing in-network settlement and using liquidity mandates/CLC Pool inventories to reduce settlement latency.

Break-even and any non-zero sCLC fee-access budgets must be evaluated on F\_cash, not gross F.

\subsection*{Break-even \& Self-Sustaining Scenarios (illustrative, update quarterly):}
\addcontentsline{toc}{subsection}{Break-even \& Self-Sustaining Scenarios (illustrative, update quarterly):}

Publish a 3-row table each quarter:

\begin{itemize}[leftmargin=*]
\item Conservative: D\_tot, V\_settle, $\tau$ $\to$ F; compare to Core Ops budget B\_core.
\item Base: same.
\item Expansion: same.
\end{itemize}
``Operational break-even'' is when F\_cash (cash-usable fee revenue) covers (Insurance top-ups + B\_core + required liquidity mandates) for 3 consecutive months under conservative assumptions. If $\chi$ is low (many fees arrive as non-convertible vouchers), break-even requires proportionally higher routed value and/or explicit conversion/subsidy policies. Therefore sCLC fee-access budgets (F\_epoch) may remain zero for extended periods until safety and operating targets are sustainably met.

As pools federate:

\begin{itemize}[leftmargin=*]
\item D(tot) tends to grow (more participants, more commitments), and
\item V(settlement) tends to rise (better routing, more netting, faster fulfillment).
\end{itemize}
Both forces push fee volume F upward.

Downstream of fees: Waterfall $\to$ policy budgets $\to$ optional sCLC budget-exit

Higher F increases the resources available to the Waterfall (insurance targets, core ops, liquidity mandates).

Only after safety and operations priorities are satisfied (insurance targets, core ops, liquidity mandates), the protocol may publish an epoch fee-credit budget F\_epoch that bounds sCLC budget-exit windows/caps into designated fee-holding pools. This makes the post-waterfall budget contestable: stakers can directly reallocate a bounded portion of pooled fee assets by exercising sCLC (e.g., injecting liquidity into specific pools), providing a ``vote with your feet'' accountability mechanism. sCLC is downstream of real settlement throughput (fulfilled commitments), not speculation, and F\_epoch may be set to zero.

Fee Flow (summary): Gross fees (in many assets)

$\to$ Waterfall (1) Insurance targets $\to$ (2) Core ops $\to$ (3) Liquidity/off-ramp mandates

$\to$ Optional: publish capped fee-credit budget F\_epoch (may be zero)

$\to$ sCLC ``budget-exit'' lets stakers direct a bounded portion of fee assets

\subsection*{5.4. Fee Pooling}
\addcontentsline{toc}{subsection}{5.4. Fee Pooling}

Liquidity providers (LPs) stake assets/vouchers into pools so that swaps and settlements can clear smoothly. They take on inventory and routing risk; service fees are the natural way to pay them.

\subsection*{LP risks \& protections (plain language)}
\addcontentsline{toc}{subsection}{LP risks \& protections (plain language)}

\begin{enumerate}[leftmargin=*]
\item Risks:
\begin{enumerate}[leftmargin=*]
\item Inventory risk: you may hold assets/vouchers that are slower to redeem or rebalance.
\item Convertibility risk: some fees arrive as non-cash vouchers; cash-usable revenue depends on $\chi$.
\item Incident risk: in extreme failures, remedies follow the disclosed loss waterfall (responsible issuer or guarantor $\to$ reserves $\to$ optional insurance $\to$ any lawfully authorized, capped reduction to optional coverage or Pool settlement claims).
\item Governance/lock risk: participation may require lockups; changes are timelocked.
\end{enumerate}
\item Protections:
\begin{enumerate}[leftmargin=*]
\item Configured limits can slow certain drains; disclosed reserves may absorb specified losses if they exist and are available. Neither eliminates loss.
\item Receipts + dashboards make issuer performance and incidents visible.
\item Policy-gated fee-access is downstream of safety/ops and may be zero---preventing ``promised yield'' dynamics.
\item Credible exit/forkability: communities can re-point/fork if governance is captured (see \S{}11.5).
\end{enumerate}
\end{enumerate}
Let:

\begin{enumerate}[leftmargin=*]
\item $\phi$ = fraction of total fees allocated to LPs (the rest can fund software, governance, guarantees, etc.)
\item K = total value of liquidity staked by LPs into the network
\end{enumerate}
The LP fee-access ex-post metric (measured from realized settlement fees; not promised returns) per period is roughly:

\[
\mathrm{FeeFlow}_{LP} \approx \frac{\phi \cdot F}{K} = \frac{\phi \cdot \tau \cdot V_{\mathrm{settle}} \cdot D_{\mathrm{tot}}}{K}
\]

This formula makes the incentive structure explicit:

\begin{enumerate}[leftmargin=*]
\item Higher settlement velocity V(settle) $\to$ more routed value $\to$ more fees $\to$ higher fee pooling.
\item More outstanding productive commitments D(tot) (claims on real output, not speculation) $\to$ a larger potential base for fee-generating fulfillment.
\item Reasonable fee rate $\tau$ and LP share $\phi$ sustain both the infrastructure and the risk-takers.
\end{enumerate}
As the network scales, LP programs may receive policy-gated fee access or sCLC allocations based on how well the system coordinates and settles real obligations, not on how much it speculates.



\section*{6. The Missing Piece: Network-Level Liquidity \& Governance}
\addcontentsline{toc}{section}{6. The Missing Piece: Network-Level Liquidity \& Governance}

Experience from the historic Sarafu Network service motivated two areas for further CLC design:

\begin{enumerate}[leftmargin=*]
\item a mechanism for LPs to \textbf{inject liquidity across pools and access sCLC}, and
\item explicit governance options as the network scales.
\end{enumerate}
One CLC design addresses both by introducing a \textbf{network clearing house (the CLC Pool)} and a \textbf{governance token (CLC)}. Other deployments may use different institutions, governance bodies, or fee models while retaining the same Commitment Pool registry and routing pattern.

\subsection*{6.1 Participation Mechanics: Downside Controls and Optionality}
\addcontentsline{toc}{subsection}{6.1 Participation Mechanics: Downside Controls and Optionality}

\textbf{Downside controls} can operate at several levels. Where configured, Swap limits and inventory checks constrain execution and on-chain transactions create persistent records. These controls reduce particular risks but do not prevent issuer failure, loss of value, contract failure, key loss, or all other losses.

A deployment may add circuit breakers, timelocks, reserves, guarantees, or an insurance runbook with a transparent loss waterfall. None is automatic. Any protection applies only if it is implemented, sufficiently funded, and published with responsible parties, coverage, caps, exclusions, claim procedures, and applicable law.

In the proposed model, governance enforcement should be narrow and reviewable: registry actions (list, suspend, or delist) should be timelocked where possible and include notice-to-cure and appeal paths, while emergency actions should require incident reporting and automatic review or sunset.

\textbf{Upside Pull (why participation may be attractive):} CLC can allow real-production and service Vouchers to function as collateral or repayment instruments in a separately documented credit facility. End-user fulfillment affects that credit only when the facility's express transaction terms define how redemption is valued and credited. Ordinary Swaps remain exchanges governed by displayed parameters and do not become loans or repayments because of asset direction.

Confederation and routing multiply fulfillment paths across pools (more routes, more netting surfaces), implementing a better matching mechanism - raising settlement velocity and fee volume - so liquidity providers and operators benefit from real settlement activity rather than speculative churn.

A compatible deployment should publish eligibility, curation, limits, health policies, refusal reasons, and transaction records so users can understand decisions. Pool Stewards and registry operators may still make admission, suspension, and routing decisions under their published governance and applicable law.



\section*{7. CLC Stewardship and the CLC Token}
\addcontentsline{toc}{section}{7. CLC Stewardship and the CLC Token}

\textbf{Design status:} This chapter is a proposed governance and liquidity-coordination design. Publication here does not mean that the CLC, stCLC, or sCLC tokens; an Insurance Fund; Waterfall contracts; gauges; mandates; fee-credit windows; or public-market programs are deployed or available in the current PWA. Each component requires a separate implementation and published policy before it applies.

\subsection*{7.1 Purpose}
\addcontentsline{toc}{subsection}{7.1 Purpose}

The proposed CLC governance layer would:

\begin{itemize}[leftmargin=*]
\item Coordinate governance for CPP networks,
\item Allocate liquidity across pools,
\item Coordinate an optional, expressly scoped and funded insurance program,
\item Maintain core infrastructure and registries,
\item Recognize the risk of liquidity providers,
\item Preserve decentralization and auditability as the network scales.
\end{itemize}
\subsection*{7.2 CLC Token Overview}
\addcontentsline{toc}{subsection}{7.2 CLC Token Overview}

In this design, CLC would be the base governance asset. Locking (staking/escrowing) CLC would mint stCLC (vote-escrow governance power) and could qualify the holder for epoch-scoped sCLC under policy.

\begin{enumerate}[leftmargin=*]
\item stCLC - a non-transferable vote-escrow receipt that represents governance voting power (and can be delegated).
\item sCLC - an epoch-scoped authorization / incentives token minted each epoch under policy and allocated to (a) stCLC holders (fee-credit authorization after the Safety Waterfall) and/or (b) approved gauges (incentives to productive liquidity and routing operators). Either allocation may be set to zero in any epoch.
\end{enumerate}
Neither stCLC nor sCLC is equity, a dividend instrument, or a guaranteed return. Policy may set sCLC issuance and/or fee-credit access to zero in any epoch.

In this proposed design, CLC would not be a community Voucher or a general medium of exchange; it would coordinate governance and policy-gated access to network resources.

\textbf{Governance Lockups (Anti-Capture; vote-escrow).} If implemented, voting power would be represented by stCLC, minted only when CLC is locked under a minimum lock period and exit cooldown set by the adopted governance parameters. Spot-held CLC would not vote. This design is intended to make hostile takeovers slower, visible, and contestable.

\textbf{sCLC Properties (Anti-Speculation).} If implemented, sCLC would be epoch-scoped, expiring or being burned at epoch end, and would function as an authorization or incentive token under caps rather than a tradable claim on profits. An adopted epoch policy could provide for emissions to approved gauges or voter incentives and could set issuance to zero in any epoch.

\textbf{Proposed Total CLC Supply:} 500,000,000 - to be minted to a CLC Vault under the custody and governance arrangements adopted for a future deployment.

A future launch could place the CLC Vault under a multisig with published signers and a rotation policy, followed by a separately approved transition to governance-controlled timelocked contracts after stated hardening milestones are met, such as two independent audits, monitoring, incident runbooks, and tested pause and fork procedures. Any use of public trading venues would require a separate safety and governance decision.

\subsection*{CLC Allocations:}
\addcontentsline{toc}{subsection}{CLC Allocations:}

\begin{itemize}[leftmargin=*]
\item \textbf{15\% Grassroots Economics Foundation (GEF):} The illustrative allocation would be permanently staked, non-transferable, and bound to the GEF multisig. It could receive governance voting and Swap-window access under adopted policy, while the underlying CLC could not be unstaked.
\item \textbf{15\% Core Team \& Early Partners:} The illustrative allocation would be staked during vesting and non-transferable until vesting ends, using a policy-set cliff and linear vesting over 24 months. Any voting during vesting and transfer restrictions would be defined in the adopted policy.
\item \textbf{30\% Endowments (Private):} The illustrative allocation would recognize early LPs providing network liquidity and use policy-set 24-month vesting.
\item \textbf{40\% Public Liquidity (DEX venue):} The illustrative allocation would be unvested and reserved for public endowment-liquidity bootstrapping.
\begin{itemize}[leftmargin=*]
\item The venue-agnostic reserve would be held in a timelocked vault and released in tranches.
\item Active deployment would be capped at 10\% of total supply across all venues at a time.
\item Each deployment would sunset after 90 days unless renewed through the adopted governance process.
\item LP positions would be governance-controlled, LP tokens would be timelocked, and public venues would remain optional.
\item CLC used for liquidity would not vote unless staked under the same lockup rules as other voters.
\end{itemize}
\end{itemize}
(These are illustrative parameters. Any adopted parameters, controls, and change process would need to be implemented and published separately.)

\subsection*{7.2.2 CLC Availability Stages}
\addcontentsline{toc}{subsection}{7.2.2 CLC Availability Stages}

\textbf{Endowment Contribution Tiers (reference valuation):} A future program could accept early endowments in staged tiers using a published reference valuation for intake and budgeting purposes. Any such reference valuation would be a disclosed governance parameter, not a promise of market price or future appreciation. Public liquidity, if separately approved on third-party venues, would be for accessibility and discovery; adopted policy should not target a price and should disclose when liquidity may be added or withdrawn.

\textbf{Endowment Covenant (Seeder Responsibility):} Under a future program, contributions would be treated as stewarded endowments intended to increase settlement capacity rather than promises of yield. Adopted policy could cap the voting influence of large endowments through defined mechanisms. Every endowment deployment should publish its purpose, expected network benefit, risks, and exit conditions.

\subsection*{7.2.3 Impact Seeding Program (CLC Eligibility for Seeding Commitment Pools)}
\addcontentsline{toc}{subsection}{7.2.3 Impact Seeding Program (CLC Eligibility for Seeding Commitment Pools)}

Published policy may allocate portions of the CLC Vault to recognize contributors who seed liquidity directly into designated Commitment Pools when that liquidity measurably increases network settlement (fulfilled redemptions), not speculative churn.

\subsection*{Eligibility (example policy, finalized by governance):}
\addcontentsline{toc}{subsection}{Eligibility (example policy, finalized by governance):}

\begin{enumerate}[leftmargin=*]
\item Seed into an approved Pool or set of Pools for a minimum duration under a rolling lockup.
\item Meet the adopted definition of ``productive'' liquidity, measured by transaction records showing support for routed Swaps that culminate in redemption or settlement within published service targets.
\item Use marginal settlement contribution, rather than TVL alone, as the basis for any reward calculation.
\end{enumerate}
An adopted program could express approved Pools as gauges so stCLC voters could transparently direct incentives toward productive settlement capacity rather than TVL.

\subsection*{Anti-gaming rules:}
\addcontentsline{toc}{subsection}{Anti-gaming rules:}

\begin{itemize}[leftmargin=*]
\item Exclude self-wash loops (same beneficial owner cycling value) and routes flagged by the router deny-list.
\item Apply per-entity caps and diminishing returns to reduce whale capture.
\item Use an observation window and delayed finalization (timelocked) to allow dispute/appeal of manipulated metrics.
\end{itemize}
\subsection*{7.3 Vote-Escrow (stCLC) + Epoch Incentives (sCLC) + Pooled Fees}
\addcontentsline{toc}{subsection}{7.3 Vote-Escrow (stCLC) + Epoch Incentives (sCLC) + Pooled Fees}

Under this proposed model, locking CLC would mint stCLC voting power and enable participation in epoch incentive decisions. During each epoch, stCLC holders would vote on gauges for approved Pools or mandates that would direct any enabled sCLC incentives to eligible liquidity and routing operators.

Separately, an implemented Safety Waterfall could fund adopted Insurance, Core Operations, and Liquidity Mandate budgets before an applicable policy defines a fee-credit budget (F\_epoch). When enabled, that policy could allow sCLC to provide capped Swap access from designated fee-holding Pools under published windows and inventory constraints.

Design rationale (``vote with your feet''): if implemented, sCLC would make post-Waterfall fee budgets contestable. Under the adopted policy, eligible stakers could reallocate a bounded portion of pooled fee assets by exercising fee-credit for allowlisted uses such as adding Pool liquidity, supporting local Voucher inventories, or purchasing coverage collateral. This would be an accountability mechanism, not a promise of yield.

\subsection*{stCLC Gauge Voting (Directing sCLC Incentives)}
\addcontentsline{toc}{subsection}{stCLC Gauge Voting (Directing sCLC Incentives)}

To reduce discretionary allocation and keep incentives tied to real settlement, the proposed design would use a gauge system containing a curated list of eligible Pools or mandates.

Each epoch:

\begin{enumerate}[leftmargin=*]
\item stCLC holders would vote on gauges for eligible Pools, portfolios, or routing mandates.
\item The implemented protocol would compute vote weights per gauge using the adopted caps and concentration controls.
\end{enumerate}
If enabled, the protocol could mint a bounded amount of sCLC incentives and distribute them to eligible liquidity providers and routing operators according to the adopted gauge and settlement-measurement policy.

\textbf{Key difference from speculation-driven AMMs:} votes do not target token price or ``APY.'' They target settlement capacity (inventory availability, routing reliability, off-ramps), measured by receipts and fulfillment outcomes.

\textbf{Waterfall Usage of Fees (policy-bound).} Under the proposed Waterfall, eligible fees would first fund any adopted Insurance Reserve Targets and Core Operations budget, followed by Liquidity Mandates. Only afterward, and only if enabled for that epoch, could the protocol publish a fee-credit budget F\_epoch that bounds sCLC budget-exit Swap access. Values, caps, and windows would be published in advance and could be tightened or set to zero under the adopted incident policy. See Section 7.4.

\subsection*{Why stake CLC?}
\addcontentsline{toc}{subsection}{Why stake CLC?}

In this design, staking or escrowing CLC would let participants direct network policy. An adopted policy could also make participants eligible for epoch-scoped sCLC and capped, epoch-bound access to a post-Waterfall fee budget. This could let eligible stakers reallocate a bounded portion of fee assets to allowlisted uses rather than relying solely on proposals and committees. It would be access to a governed resource under disclosed caps, not a claim on profits or dividends.

\subsection*{Service-Fee Enforcement Option}
\addcontentsline{toc}{subsection}{Service-Fee Enforcement Option}

(a) \textbf{Factory Gating.} A registry operator may deploy CPP pools via a \textbf{PoolFactory} that wires a \textbf{FeeHook} into the Vault/Fee Registry; fees can auto-route to a configured fee recipient via the Waterfall per policy. Pools missing the FeeHook may be ineligible for that registry.

(b) \textbf{Registry Gating.} A registry profile may make only registered pools discoverable by its routers and SDKs. Independent registries can define their own listing and fee rules.

(c) \textbf{Router Policy.} Routers may \textbf{refuse routes} that touch unregistered pools or pools with invalid FeeHooks for the selected profile. SDK invariant checks can enforce this.

(d) \textbf{Programmatic Attestations.} Liquidity mandates and LP programs may require FeeHook compliance; non-compliant pools may lose routing and liquidity support within that profile.

\textbf{Result:} ``Mandatory'' is profile-specific: a non-paying pool may be unroutable in one registry or service profile while remaining free to operate locally or join another registry.

\textbf{Fork/Exit Note.} This enforcement applies only to the selected discovery profile. Pools remain free to operate outside a given registry, and independent routers/registries may exist. This preserves credible exit if governance is captured: a fork can deploy alternative registries/routers and pools can re-register there without changing the underlying CPP primitives.

Compatible clients should permit Users and Pools to select available alternative profiles or registries. Registry enforcement should not disable otherwise functional local Pools, and applicable fee policies should be displayed before a Swap or contribution.

\subsection*{sCLC Emission \& Budget-Exit Windows}
\addcontentsline{toc}{subsection}{sCLC Emission \& Budget-Exit Windows}

For an implemented epoch, after the Waterfall funds the adopted Insurance, Core Operations, and Liquidity Mandate budgets, the protocol could publish a fee-credit budget F\_epoch whose value may be zero. Under the proposed policy, sCLC would confer a pro-rata User limit based on stCLC voting power:

\[
\mathrm{limit}_{\mathrm{user,epoch}} = F_{\mathrm{epoch}} \times \frac{\mathrm{stCLC}_{\mathrm{user}}}{\mathrm{stCLC}_{\mathrm{total}}}
\]

Within published windows and caps, and subject to inventory, an implemented policy could permit sCLC to be exercised to Swap fee assets out of designated fee-holding Pools and into listed assets or Vouchers. This would be a bounded budget-exit mechanism rather than passive income; governance could set F\_epoch to zero or tighten, pause, or geofence access under the adopted compliance and incident policy.

\textbf{DEX Float Reduction (optional, non-speculative):} If measured CLC DEX float exceeded an adopted policy cap, published policy could authorize a capped, TWAP-limited repurchase solely to reduce external float and governance-attack surface. Acquired CLC would be retired or placed in a disclosed non-voting sink. Such a program should have no price target, could be set to zero, and should halt under its published incident and reserve rules.

\subsection*{Guardrails (policy parameters; on-chain):}
\addcontentsline{toc}{subsection}{Guardrails (policy parameters; on-chain):}

\begin{itemize}[leftmargin=*]
\item \textbf{Trigger-based:} only if DEX float > \textbf{X\%} for \textbf{Y} days
\item \textbf{Cap:} max \textbf{Z\%} of monthly fees and/or max \textbf{W\%} of DEX daily volume
\item \textbf{Execution:}\textbf{TWAP + random time delays}+ \textbf{no announcement of exact timing}
\item \textbf{Emergency stop:} automatic stop if \textbf{insurance ratio < threshold}
\item \textbf{Disclosure:}``\textbf{Not intended to support price; settlement does not depend on DEX price}''
\end{itemize}
\textbf{Treasury Liquidity Cache (non-distributive):} An adopted policy could use a disclosed portion of eligible fees to maintain protocol-controlled liquidity positions for stated network functions, such as off-ramp buffers or rebalancing inventory, under published caps and controls.

\textbf{7.3.1 Portfolio Pools:} Direct Seeding, Voted Allocations, and sCLC-Directed Liquidity (Examples)

Commitment Pools can be curated as ``portfolio pools'': pools that list redeemable commitments aligned to a mission domain (e.g., ecosystem support services, humanitarian support, health \& wellness). Portfolio pools make it easy for liquidity providers to target real-world outcomes without requiring a single central issuer.

\subsection*{There are three complementary ways to support a specific portfolio pool:}
\addcontentsline{toc}{subsection}{There are three complementary ways to support a specific portfolio pool:}

A) Seed directly into the pool: deposit accepted assets/vouchers into the pool's Vault, increasing inventory and routing capacity (subject to listings, limits, and reserve policy).

B) Vote allocations into the pool: propose and approve Liquidity Mandates (Waterfall allocations) that seed or backstop designated portfolio pools with time-bounded mandates and sunset/review.

C) Stake CLC and direct sCLC swaps: when enabled for an epoch, stakers can exercise sCLC budget-exit swaps to reallocate a bounded portion of post-waterfall fee assets into specific portfolio pools (e.g., by swapping fee assets into the pool's inventories or accepted liquidity assets), strengthening the pools they believe most improve settlement and mission outcomes. This is an accountability mechanism under caps/windows---not a claim on profits.

Note: Portfolio pools remain sovereign. They can be canonical (discoverable via official registries/routers) or independent (discoverable via independent registries/routers), without changing the underlying CPP primitives.

\subsection*{7.3.2 Curating Portfolio Pools (Including Certifications)}
\addcontentsline{toc}{subsection}{7.3.2 Curating Portfolio Pools (Including Certifications)}

Any lawful Steward, such as an individual, cooperative, community group, multisig, service operator, or public agency, could curate a portfolio Pool by defining a listing policy, publishing a Value Index method, configuring limiters, and requiring clear redemption evidence and fallback remedies. Portfolio Pools could be specialized or mixed.

Certifications could be used to support trust and risk assessment, but should be modeled as attestations that affect eligibility and risk treatment rather than as profit tokens. Two possible patterns are:

A) Attestation Certificates (non-transferable or registry-bound): a verifier issues an attestation that a Voucher issuer or project meets stated criteria (methodology, safeguards, monitoring). A Pool could use attestations to whitelist listings, apply disclosed valuation discounts, widen or narrow limits, or qualify for an optional insurance program.

B) Audit/Verification Service Vouchers (redeemable commitments): a token represents a redeemable verification service (who will verify what, by when, under what standard). Pools/projects can purchase these vouchers to fund monitoring and strengthen integrity.

In both cases, the economic claim remains the underlying redeemable commitment; certifications modify risk and eligibility rather than creating entitlement to fees, profits, or residual assets.

\subsection*{7.4 Waterfall Policy \& Budgets}
\addcontentsline{toc}{subsection}{7.4 Waterfall Policy \& Budgets}

In the proposed model, eligible fee inflows from Pool usage, routing, or a network rake would be allocated by an implemented Waterfall under the adopted governance process.

\textbf{Fee Asset Eligibility \& Conversion (cash vs. in-kind).} A future deployment could receive fees in mixed assets. Its Waterfall would distinguish:

(i) Cash-eligible fee assets (E\_cash): allowlisted stablecoins/cash-equivalents and (optionally) major liquid tokens that may be converted to fund fiat-denominated insurance payouts and core operating costs; and

(ii) In-kind fee assets (E\_kind): non-fiat-redeemable vouchers and other non-convertible assets that may be redeployed for in-network settlement support, local mandates, or in-kind operating needs, but do not count toward fiat-denominated insurance/ops obligations.

A deployment using this model should publish a Conversion Policy for fungible assets covering allowlists, caps, slippage limits, TWAP windows, and reporting. Voucher pricing would remain governed by Pool Value Indices and Swap Limiters; treasury conversion would not determine Voucher valuation.

\subsection*{Waterfall Priorities:}
\addcontentsline{toc}{subsection}{Waterfall Priorities:}

\begin{enumerate}[leftmargin=*]
\item \textbf{Insurance Reserve Target}-- Fund to a policy target (risk-weighted by pool class, fulfillment rate, issuer concentration, and limit utilization).
\item \textbf{Core Operations}-- Legal, advocacy, IEC, infra, audits, observability.
\item \textbf{Liquidity Mandates}-- Endowments into target pools/routers to improve settlement velocity, including (optionally) interoperability mandates: bridge/adaptor maintenance, confederation routing pilots, and cross-network liquidity backstops under published caps and sunset reviews.
\item \textbf{DEX Float Reduction (optional, non-speculative):} If measured CLC DEX float exceeded an adopted policy cap, published policy could authorize a capped, TWAP-limited repurchase solely to reduce external float and governance-attack surface. Acquired CLC would be retired or placed in a disclosed non-voting sink. Such a program should have no price target, could be set to zero, and should halt under its published incident and reserve rules.
\item \textbf{CLC Pool Fee-Access Budget:} Allocate remaining eligible fee assets to the CLC Pool (cash-eligible E\_cash by default; E\_kind only if explicitly allowlisted per program) and publish the epoch fee-access budget F\_epoch (may be zero), which bounds sCLC swap-access windows/caps as defined in \S{}7.4.
\end{enumerate}
\textbf{KPI-Linked Budgets.} A deployment could use advisory Pool-health data to inform Waterfall parameters, including fulfillment rate, reserve adequacy, limit utilization, routing results, guarantor performance, and redemption latency. Any adopted edit process should disclose its authorization, delay, and recordkeeping rules.

A possible contributor flow under this proposal would be:

\begin{enumerate}[leftmargin=*]
\item Liquidity Providers would contribute eligible assets to the CLC Vault under separately adopted endowment terms and could receive CLC tokens as governance participants.
\item CLC holders could stake under the adopted rules to receive stCLC voting rights and, where enabled, sCLC under step 4.
\item Participating Pools could send a disclosed percentage of their fees to an implemented Waterfall contract, which would allocate eligible fees to:
\begin{enumerate}[leftmargin=*]
\item an Insurance and Operations Vault funded primarily by cash-eligible fee assets or policy-authorized conversions, with in-kind fees excluded from fiat-denominated obligations;
\item participating CPs according to adopted voting rules;
\item an optional DEX Float Reduction program only when its trigger conditions are met, otherwise zero; and
\item a CLC Fee Budget Vault holding post-Waterfall eligible fee assets, where an enabled sCLC policy could provide capped, epoch-bound Swap access to allowlisted assets or programs.
\end{enumerate}
\item After the Safety Waterfall, an adopted epoch policy could mint sCLC and allocate it to eligible stCLC holders as fee-credit authorization, approved gauges as incentives, or both; either allocation could be zero.
\item A CLC Fee Budget Vault would hold post-Waterfall pooled fee assets. When enabled, an adopted sCLC policy could provide capped, epoch-bound authority to execute allowlisted deployments from designated fee-holding vaults within published windows and caps. This would not be a claim on Vault ownership.
\end{enumerate}


\section*{8. Technical Scope \& Growth}
\addcontentsline{toc}{section}{8. Technical Scope \& Growth}

This section describes optional and future technical work. It is not a list of features guaranteed to be present in the current public PWA.

Priorities:

\begin{itemize}[leftmargin=*]
\item \textbf{Routing protocols across} Cosmo-Local Credit and other compatible Pools; SDKs and index/limit discovery APIs.
\item \textbf{Bridges to external DEX} registries; Time-locked contract/escrow for cross-domain settlement.
\item Support for the long tail of \textbf{micro-pools} (including personal pools UX).
\item \textbf{Auditable registries} for vouchers, pools, limits, values, and fee policies.
\item \textbf{Fiat/stable on/off-ramp connectors}: a partner/rail registry (by jurisdiction), checkout/invoice flows, and UI ``network profiles'' that can geofence features and require attestations for cash-equivalent redemptions.
\item \textbf{Bridges to external liquidity venues} (DEXs and other registries) for fungible assets (e.g., stable cash-equivalents and major liquid tokens), using time-locked escrow/HTLC where cross-domain atomicity is unavailable. Purpose: rebalancing, on/off-ramp liquidity, and risk-managed settlement support - not pricing of redeemable vouchers by speculative curves.
\item \textbf{Treasury conversion \& settlement acceleration}: policy-capped conversion of cash-eligible fee assets (E\_cash) held in the CLC Pool into stables/fiat when needed for insurance/ops and to maintain off-ramp liquidity---while prioritizing in-network settlement and using liquidity mandates/rebalancing to reduce settlement latency.
\end{itemize}
\textbf{DEX Interop Boundary:} DEX adapters are for fungible liquidity management (stables, rebalancing, exit ramps), not for defining the value index of redeemable commitments. Voucher pricing remains governed by each pool's Value Index + limits + inventories; DEX prices may be used only as an auxiliary reference for fungible assets and must be guarded against manipulation (caps, TWAP/medianization, deny-lists, and incident pauses).

\subsection*{8.1 Router \& SDK Norms}
\addcontentsline{toc}{subsection}{8.1 Router \& SDK Norms}

\textbf{Public Discovery.} Routers must query public registries of voucher listings, value indices, limits, fees, and inventories; cache with freshness bounds.

\textbf{Multi-Profile Discovery (Confederation).} Routers/SDKs may support multiple registry roots (``network profiles'') and must surface to users/pools which profile a route uses (registry root, policy constraints, bridge adapters). Cross-profile routes must satisfy the strictest applicable caps/escrow requirements and must be auditable hop-by-hop (quote $\to$ receipt mapping).

\textbf{Path Policies.} Deny-list toxic routes (known bad bridges/pools) and enforce per-route caps and minimum health scores (reserve adequacy, SLA adherence).

\textbf{Fees \& Caps.} Routing fees expressed per-hop; routers may add a small discoverability fee within policy bounds; hard caps apply under stress (utilization spikes).

\textbf{Atomicity \& Escrow.} Prefer atomic multi-hop where possible; otherwise use HTLC/escrow with conservative timeouts and explicit abort paths.

\textbf{Batch Netting \& Rebalancing.} Routers may also perform batch netting runs across opted-in pools by collecting rebalance intents and searching for multilateral cycles/chains that satisfy each pool's constraints. Norms:

(i) publish a machine-readable ``rebalance receipt'' summary (cycles executed, total off-set value, fees charged),

(ii) enforce conservative per-epoch caps and health-score gating,

(iii) reject any route that violates a pool's allow/deny policies or exceeds limiter windows/caps,

(iv) keep batch execution auditable (deterministic inputs $\to$ receipts) to support dispute resolution.

\textbf{SDK Guarantees.} Provide (i) deterministic quote $\to$ receipt mapping, (ii) invariant checks per hop, (iii) human-readable failure codes, (iv) audit-friendly logs.

\subsection*{8.1.1 Minimum Confederation Compatibility Contract (MCC)}
\addcontentsline{toc}{subsection}{8.1.1 Minimum Confederation Compatibility Contract (MCC)}

To be routable across profiles (and thus across confederated networks), a pool ecosystem must publish the following in a machine-readable way:

\begin{enumerate}[leftmargin=*]
\item Registry roots: canonical addresses for voucher registry, pool registry, value index registry, limiter registry, and fee policy registry (or a single root that deterministically resolves these).
\item Receipt standard: every hop must emit/return a receipt that references (a) registry root/profile used, (b) voucher/token in/out, (c) value index version or timestamp, (d) limiter window/cap snapshot, (e) fees charged, and (f) inventory check result.
\item Health endpoints: per-pool signals required for routing policies---reserve adequacy, SLA adherence, limiter utilization, and incident state---plus freshness bounds.
\item Policy constraints: explicit allow/deny policies (bridges, counterparties, voucher classes) and required escrow/HTLC requirements for non-atomic hops.
\item Failure codes: human-readable, deterministic failure codes so clients can explain why a route was refused (limits, inventory, policy, escrow, incident pause).
\end{enumerate}
Networks may implement additional features (insurance overlays, compliance hooks, arbitration modules), but these must remain profile-scoped and must not be required for basic CPP compatibility.

\subsection*{8.2 Licensing \& Transparency}
\addcontentsline{toc}{subsection}{8.2 Licensing \& Transparency}

The protocol contracts are \textbf{EVM-compatible}. Contracts under the protocol's \texttt{src} directory are published under \textbf{AGPL-3.0} except for unmodified third-party components, such as identified Solady contracts, that retain their original terms. Published source, ABIs, and deployment instructions support independent verification. Each deployment must separately disclose its code version, build provenance, canonical addresses, administrator and upgrade powers, audit status, registry mirrors, and any timelock protections; none of those protections should be assumed merely from protocol compatibility.

\textbf{Fork Kit (Proposed Deliverable).} A mature deployment should maintain a ``fork kit'' that includes:

(i) deterministic deployment scripts; (ii) registry snapshot/export tooling; (iii) a documented procedure to re-point routers/SDKs to a new registry root; and (iv) a pool steward checklist for exiting canonical registries safely (including fee-hook redirection options where supported).

\subsection*{Example: Minimum Exit Checklist (publish + test annually):}
\addcontentsline{toc}{subsection}{Example: Minimum Exit Checklist (publish + test annually):}

\begin{enumerate}[leftmargin=*]
\item How to export registry snapshots + receipts.
\item How to repoint routers/SDKs to a new root.
\item How to migrate insurance scope (or explicitly terminate it).
\item How to honor outstanding vouchers during migration (notice-to-redeem + remedy options).
\end{enumerate}
\textbf{Why shared-source licensing + a Fork Kit:} Confederation rewards compatibility. Networks that fork and improve routers, registry tooling, bridge adapters, or observability can still route with CLC if they remain CPP-compatible, expanding settlement paths and strengthening the whole mesh. Improvements to AGPL-covered shared infrastructure remain shareable under its terms, while identified third-party components retain their applicable terms.



\section*{9. Economics for LPs}
\addcontentsline{toc}{section}{9. Economics for LPs}

\textbf{Design status:} The mechanisms and figures in this chapter are illustrative proposals, not current App features, promised returns, or offers. A deployment must publish its actual fees, budgets, assets, risks, and eligibility rules before participation.

\subsection*{9.1 Revenue Streams}
\addcontentsline{toc}{subsection}{9.1 Revenue Streams}

\begin{enumerate}[leftmargin=*]
\item \textbf{Network Fee Rake $\to$ Policy Pools.} Under the proposed model, a disclosed portion of participating Pool fees would be routed through an implemented Waterfall to adopted insurance, Core Operations, and Liquidity Mandate budgets, followed only where enabled by \textbf{sCLC Swap windows}.
\item \textbf{Policy-gated fee-credit (ex-post)}
\item \textbf{Routing Fees}: Fees from multi-hop routes discovered by routers.
\begin{enumerate}[leftmargin=*]
\item Rebalancing / Netting Fees (optional): Fees earned for executing batch netting cycles and inventory rebalancing routes that reduce imbalance and increase successful settlement throughput (ex-post, policy-bound).
\end{enumerate}
\item \textbf{Curation \& Validation Fees:} pools may allocate a disclosed portion of fees to listing/verification/monitoring roles (curators, auditors, claims modules) under published mandates.
\end{enumerate}
\subsection*{9.2 Illustrative Fee Math (Example Only)}
\addcontentsline{toc}{subsection}{9.2 Illustrative Fee Math (Example Only)}

\begin{itemize}[leftmargin=*]
\item Per-pool usage fee: 30--500 bps depending on voucher class and risk tier.
\item Network rake: 10--30\% of per-pool fees -> Waterfall -> CLC Pool.
\begin{itemize}[leftmargin=*]
\item \textbf{Worked example} ( ``2\% fee'' case). If a pool charges f\_p = 2.00\% and the network rake share is r\_p = 20\% of that pool's fees, then the effective network fee rate on routed value is:
\item $\tau$\_p = f\_p $\cdot$ r\_p = 2.00\% $\cdot$ 20\% = 0.40\% = 40 bps.
\item Convertibility matters. If only $\chi$ = 25\% of fee inflows are cash-eligible/convertible (E\_cash), then cash-usable effective revenue is \textasciitilde{}10 bps (40 bps $\times$ 0.25). Therefore, meaningful sCLC fee-access budgets (F\_epoch) require both high settlement throughput and sufficient $\chi$; otherwise F\_epoch may remain zero for long periods.
\end{itemize}
\item Routing fee: 5--20 bps across hops.
\item Distribution to Waterfall (policy-bound)
\item \textbf{Net LP Credit Access} drivers: local Swap usage, routing volume, policy-deployed liquidity access, less losses and any authorized reductions to optional coverage or Pool settlement claims.
\end{itemize}
*\textbf{Accounting only: Any annualized figures are ex-post metrics of policy-gated Swap access to pooled fees (not promised returns) and may be zero or negative after losses or authorized coverage reductions.}*

\subsection*{Downside Examples (ex-post):}
\addcontentsline{toc}{subsection}{Downside Examples (ex-post):}

\begin{itemize}[leftmargin=*]
\item Inventory Loss Case: losses from defaults or redemption delays reduce fee-credit access by X.
\item Run-Protection Case: limiter-triggered throttling reduces settlement flow, lowering F temporarily.
\item Policy Case: governance sets fee-credit budget F\_epoch = 0 (no sCLC exit) during incidents or rebuild phases.
\end{itemize}
\emph{Any numerical schedule would require separate adoption, disclosure, and enforcement under the applicable deployment's governance and terms.}



\section*{10. Comprehensive Risk Framework}
\addcontentsline{toc}{section}{10. Comprehensive Risk Framework}

We classify risk into \textbf{ten} categories. For each we list \emph{Threats}, \emph{Indicators}, and possible \emph{Controls} (Preventive/Detective/Corrective), plus \emph{Stress Tests} and a \emph{Risk Appetite} statement. These are design recommendations, not representations that every control is deployed, effective, or sufficient. Limits, reserves, guarantees, monitoring, insurance, and governance processes cannot eliminate loss.

\subsection*{10.1 Protocol \& Smart Contract Risk}
\addcontentsline{toc}{subsection}{10.1 Protocol \& Smart Contract Risk}

\begin{itemize}[leftmargin=*]
\item \textbf{Threats}: Contract bugs, upgrade errors, misconfigured limits/fees.
\item \textbf{Indicators}: Audit findings, unexplained vault movements, reverts/spikes.
\item \textbf{Controls}:
\begin{itemize}[leftmargin=*]
\item \emph{Preventive}: Independent audits, formal verification where critical, minimal privileged roles, least-privilege vaults.
\item \emph{Detective}: On-chain monitors, invariant checks, receipt reconciliation.
\item \emph{Corrective}: Timelocked upgrades, emergency pause with public criteria, fork/migration path.
\end{itemize}
\item \textbf{Stress Tests}: Simulate paused oracles, inventory shortfalls, routed burst traffic.
\item \textbf{Risk Appetite:} Low; require audits before scaling network exposure (outstanding obligations and routed volume).
\end{itemize}
\subsection*{10.2 Economic/Market Risk (Liquidity, Runs, Oracle)}
\addcontentsline{toc}{subsection}{10.2 Economic/Market Risk (Liquidity, Runs, Oracle)}

\begin{itemize}[leftmargin=*]
\item \textbf{Threats}: Thin inventory, bank-run dynamics, oracle manipulation/latency.
\item \textbf{Indicators}: Limit utilization >80\%, widening spreads, frequent limit rejections.
\item \textbf{Controls}:
\begin{itemize}[leftmargin=*]
\item \emph{Preventive}: Swap Limiter windows/caps; tiered limits; minimum reserve ratios; medianized oracles with failover constants.
\item \emph{Detective}: Utilization \& spread dashboards; variance alarms on value index updates.
\item \emph{Corrective}: Widen margins; tighten caps; route-around policies; temporary fee surcharges to dampen flow.
\end{itemize}
\item \textbf{Stress Tests}: Price shocks ($\pm$50\%), redemption surges (5--10$\times$ baseline), oracle outages.
\item \textbf{Risk Appetite}: Moderate, bounded by policy thresholds.
\end{itemize}
\subsection*{10.3 Credit/Voucher Issuance Risk}
\addcontentsline{toc}{subsection}{10.3 Credit/Voucher Issuance Risk}

\begin{itemize}[leftmargin=*]
\item \textbf{Threats}: Over-issuance vs. capacity; issuer default; mis-specified redemption windows.
\item \textbf{Indicators}: Fulfillment rate $\downarrow$, aging vouchers $\uparrow$, guarantor exposures $\uparrow$.
\item \textbf{Controls}:
\begin{itemize}[leftmargin=*]
\item \emph{Preventive}: Issuer due diligence; pool operator guarantor/bond requirements; issuance quotas; clarity-first voucher terms.
\item \emph{Detective}: Cohort aging reports; fulfillment SLA tracking; guarantor performance logs.
\item \emph{Corrective}: Tighten issuance, require top-ups, delist/route-around, trigger insurance.
\end{itemize}
\item \textbf{Stress Tests}: Issuer insolvency; regional shock to production.
\item \textbf{Risk Appetite}: Moderate for diversified issuers; low for concentrated exposure.
\end{itemize}
\subsection*{10.4 Redemption/Operations Risk}
\addcontentsline{toc}{subsection}{10.4 Redemption/Operations Risk}

\begin{itemize}[leftmargin=*]
\item \textbf{Threats}: Inability to honor redemptions due to logistics, supply chain, or custody failures.
\item \textbf{Indicators}: Redemption latency > SLA; stockouts; ticket backlogs.
\item \textbf{Controls}: Redemption buffers; emergency backstops; off-chain attestation flows; multi-venue redemption options.
\item \textbf{Stress Tests}: 2--4$\times$ redemption spikes; facility outages.
\item \textbf{Risk Appetite}: Low; protect end-users.
\end{itemize}
\subsection*{10.5 Governance Risk}
\addcontentsline{toc}{subsection}{10.5 Governance Risk}

\begin{itemize}[leftmargin=*]
\item \textbf{Threats}: Steward capture; rushed parameter edits; conflicts of interest.
\item \textbf{Indicators}: Concentrated voting power; frequent emergency actions; policy churn.
\item \textbf{Controls}: Quorums; timelocks; delegated voting with transparency; conflict disclosures; veto/appeal mechanisms; forkability.
\item \textbf{Risk Appetite}: Low; emphasize transparency and recourse.
\item \textbf{Stress Tests}: Adversarial proposals; bribery attempts.
\begin{itemize}[leftmargin=*]
\item \textbf{Additional Stress Test: Hostile Voting Capture (Public Market Accumulation).}
\end{itemize}
\end{itemize}
Scenario: An external actor accumulates a large fraction of CLC, delegates votes to a small set of accounts, and proposes to (i) redirect Waterfall allocations, (ii) weaken listing/delisting standards, (iii) drain insurance via permissive claims, or (iv) force liquidity mandates into self-serving pools.

\begin{itemize}[leftmargin=*]
\item \textbf{Controls (must all be true):}
\begin{itemize}[leftmargin=*]
\item \textbf{Governance requires lockups} (no instant voting power from spot purchases).
\item \textbf{Timelocks} on all critical actions, with public alerts and a defined response window.
\item \textbf{Supermajority + higher quorum tiers} for Waterfall, registry root, insurance policy, and emergency powers.
\item \textbf{Transparent delegation}+ concentration monitoring triggers (automatic escalation to higher thresholds).
\item \textbf{Appeal + incident process} and an explicit \textbf{fork-and-migrate} procedure if values drift or capture occurs.
\end{itemize}
\end{itemize}
\subsection*{10.6 Legal \& Compliance Risk}
\addcontentsline{toc}{subsection}{10.6 Legal \& Compliance Risk}

\begin{itemize}[leftmargin=*]
\item \textbf{Threats}: Vouchers deemed regulated instruments; KYC/AML obligations; cross-border restrictions.
\item \textbf{Indicators}: Jurisdictional flags; regulator inquiries.
\item \textbf{Controls}: Modular compliance hooks (allow/deny lists, attestations); geofenced UIs; legal reviews per class; disclosures.
\item \textbf{Stress Tests}: Jurisdictional bans; counterparty de-listings.
\item \textbf{Risk Appetite}: Low; comply or geofence.
\end{itemize}
\subsubsection*{10.6.1 Legal Positioning \& Token Treatment (Summary)}

\begin{enumerate}[leftmargin=*]
\item \textbf{Verifiable infrastructure.} The protocol contracts are EVM-compatible and their source is published under the terms and third-party exceptions identified in the protocol repository. Publication makes independent review possible but does not mean that a contract or deployment has been audited. Each deployment should disclose its code version, build provenance, addresses, administrator powers, and audit status.
\item \textbf{Proposed token posture.} In this design, CLC would be a governance and access token. A separately implemented policy could allow staking to mint sCLC that confers policy-gated Swap rights into fee-holding Pools under published caps and windows. The design provides no dividends, profit share, or residual rights.
\item \textbf{Proposed representation policy.} Materials for any CLC, stCLC, or sCLC deployment should not promise profit or market the assets using expected-return language.
\item \textbf{Possible jurisdiction controls.} Depending on applicable law, a deployment could use: (i) geofenced interfaces or RPC access; (ii) attestation gates for restricted classes; (iii) promotion restrictions; (iv) per-Voucher-class legal review; and (v) programmatic controls that can disable sCLC Swap windows under published policy.
\item \textbf{Proposed endowment notice.} Any adopted terms should explain that access based on staked CLC may be disabled or reduced for compliance, operational, or risk reasons and whether compensation or another remedy is available. See \S{}17.3 for plain-language instrument definitions.
\end{enumerate}
\subsection*{10.7 Routing \& Cross-Domain Risk}
\addcontentsline{toc}{subsection}{10.7 Routing \& Cross-Domain Risk}

\begin{itemize}[leftmargin=*]
\item \textbf{Threats}: Partial fills, stuck hops, bridge exploits, MEV (router perspective): artificially inflate route hops, front-running value index (if timelocked).
\item \textbf{Indicators}: HTLC expiries; escrow backlogs.
\item \textbf{Controls}: Atomicity where possible; escrow/HTLC with conservative timeouts; deny risky bridges; per-route fees and caps slash router stake; implement path efficiency score (including retroactively, see OP style challenge games), capping fees.
\item \textbf{Stress Tests}: Bridge halt; chain reorgs.
\item \textbf{Risk Appetite}: Low to moderate; whitelist bridges.
\end{itemize}
\subsection*{10.8 Custody \& Key Management Risk}
\addcontentsline{toc}{subsection}{10.8 Custody \& Key Management Risk}

\begin{itemize}[leftmargin=*]
\item \textbf{Threats}: Key loss/compromise; signer collusion.
\item \textbf{Indicators}: Anomalous signer behavior; threshold changes.
\item \textbf{Controls}: Multisig/threshold schemes; hardware security; signer rotation; monitoring; withdrawal rate limits.
\item \textbf{Risk Appetite}: Low.
\end{itemize}
\subsection*{10.9 Reputation \& Social Risk}
\addcontentsline{toc}{subsection}{10.9 Reputation \& Social Risk}

\begin{itemize}[leftmargin=*]
\item \textbf{Threats}: Misaligned incentives harming communities; poor redemption experiences, insurance gamification.
\item \textbf{Indicators}: Community feedback; complaint ratios; social sentiment.
\item \textbf{Controls}: NVC-aligned evaluation; grievance redressal; transparent reporting.
\item \textbf{Risk Appetite}: Low; prioritize care and fairness.
\item \textbf{Mis-selling controls:} standardized consumer disclosures in-UI; mis-selling incident log; sanctions ladder for violators (warnings $\to$ suspensions $\to$ delistings).
\end{itemize}
\subsection*{10.10 Concentration \& Fragmentation Risk}
\addcontentsline{toc}{subsection}{10.10 Concentration \& Fragmentation Risk}

\begin{itemize}[leftmargin=*]
\item \textbf{Threats}: Dependence on a few issuers/pools; incompatible forks.
\item \textbf{Indicators}: HHI by issuer/pool; routing failures across clusters.
\item \textbf{Controls}: Diversification targets; publish registries/indices/limits; maintain routing bridges; encourage standard adherence.
\end{itemize}


\section*{11. Governance Mechanics}
\addcontentsline{toc}{section}{11. Governance Mechanics}

\textbf{Design status:} This chapter is a governance template, not a representation that the current App uses a CLC token vote, timelocks, shared insurance, a claims process, or every control described below. A deployment must identify its actual decision-makers, authorities, contracts, processes, and policies.

\begin{itemize}[leftmargin=*]
\item \textbf{Constitutional Values}: Care for People, Care for the Environment, Fairness, Reciprocity, Non-Dominance, Resilience.
\item \textbf{Proposal Types}: Fee/limit/index edits; liquidity mandates; pool listings/delistings; insurance payouts; parameter guardrails.
\item \textbf{Process}: Intake $\to$ Evaluation (template) $\to$ Risk review $\to$ On-chain vote (staked CLC) $\to$ Timelock $\to$ Execution.
\item \textbf{Quorum \& Thresholds}: Parameterized per class (e.g., higher for value-index edits and emergency pauses).
\item \textbf{Delegation}: Optional delegate system with public mandates and recall.
\item \textbf{Circuit Breakers}: Emergency pause with criteria; automatic resume conditions; post-mortems required.
\item \textbf{Transparency}: All edits/flows logged; dashboards for fulfillment, reserves, utilization, routing, guarantors.
\end{itemize}
\textbf{Registry Governance (Listing / Suspension / Delisting).} A CLC-compatible deployment may maintain discovery registries for Vouchers, Tokens, and Pools. Depending on the governance model, authorized controls may add, update, suspend, or remove (``delist'') registry entries through on-chain voting, multisig approvals, board decisions, cooperative resolutions, public mandates, or other accountable processes. Published registry rules should make status conditional and may identify repeated non-fulfillment, fraud or misrepresentation, unsafe contract behavior, or persistent violation of published principles as grounds for suspension or delisting. Routers may route around entries that the selected registry has delisted. Where feasible, the adopted process should provide notice, an opportunity to remedy, and an appeal path; emergency delisting should require a public incident report and automatic review or sunset.

\subsection*{Prohibited Listings (non-negotiable):}
\addcontentsline{toc}{subsection}{Prohibited Listings (non-negotiable):}

\begin{enumerate}[leftmargin=*]
\item Instruments that directly fund or incentivize ecological destruction beyond agreed boundaries, violence/weaponization, coercive extraction, or systemic abuse.
\item Any voucher class lacking clear redemption terms, accountability, and remedy pathways.
\end{enumerate}
Under this governance template, the prohibited list would be versioned, publicly auditable, and changeable only through the adopted critical-action threshold and timelock, illustrated as Q3 + T3 in Appendix D.

\subsection*{11.1 Index \& Limit Governance}
\addcontentsline{toc}{subsection}{11.1 Index \& Limit Governance}

\textbf{Timelocked Edits.} A deployment following this template should execute Value Index and Swap Limiter parameter edits after a public timelock. Any emergency path should use a separately disclosed authorization process and include automatic sunset or review.

\textbf{Quorum \& Thresholds.} The template proposes higher quorum or approval thresholds for (a) Value Index base changes and (b) global Limit Tier changes, intermediate thresholds for Pool-specific third-party changes, and standard thresholds for routine fee changes.

\textbf{Publish Feeds.} A participating deployment should publish, for each Pool, the on-chain index variables, oracle sources or medians, update cadence, limit windows and caps, and failure modes or safe constants.

\textbf{Emergency Pause Criteria.} A participating deployment should pre-declare conditions such as an oracle outage, high limit utilization combined with redemption failures, or an invariant failure, together with resume checks and post-incident review requirements.

\subsection*{Ex. Public Index Feed (per pool, per voucher)}
\addcontentsline{toc}{subsection}{Ex. Public Index Feed (per pool, per voucher)}

\begin{itemize}[leftmargin=*]
\item Symbol: e.g., Maize\_50kg@IssuerY
\item Reference Unit: ``Index Unit'' (IUX)
\item Valuation: 30.000 IUX
\item Source: Median(Oracles: local market survey, ministry bulletin, CLC baseline)
\item Update Cadence: daily at 18:00 EAT; Timelock: 24h
\item Failure Mode: freeze at last-good, widen limiter bands by +20\%, pause at 72h outage
\item Rationale: published notes + diff from prior update
\item Signers: multisig addresses; quorum threshold
\end{itemize}
\subsection*{11.2 Insurance Fund Runbook}
\addcontentsline{toc}{subsection}{11.2 Insurance Fund Runbook}

\textbf{Optional design only.} This runbook applies only to a deployment that has expressly adopted and funded an Insurance Fund and published the covered events, eligible claimants, responsible entity, assets, limits, exclusions, evidence, process, and governing terms. Neither the current App nor GEF provides coverage merely because this design appears in the White Paper.

\textbf{Triggers.} (i) Issuer default/non-fulfillment; (ii) Pool insolvency (reserve shortfall vs. bonds); (iii) Bridge/escrow loss impacting redeemability.

\textbf{Assessment.} Convene risk committee; reconcile receipts, vault balances, guarantor bonds, and redemption tickets; publish incident ledger.

\textbf{Illustrative Loss Waterfall.} Where each layer exists and lawfully applies: (1) responsible issuer bonds or guarantor stakes $\to$ (2) Pool-level reserves $\to$ (3) Network Insurance Fund $\to$ (4) a temporary reduction to an optional insurance payout or Pool settlement claim, but only where pre-existing applicable terms expressly authorize it, law permits it, and every required consent and process is satisfied $\to$ (5) lawful recovery for proven fraud or abuse. A coverage or settlement haircut does not reduce an Issuer's underlying Voucher commitment or alter an on-chain balance unless valid pre-existing terms and applicable law expressly permit that result and any required Holder consent is obtained.

\textbf{Coverage Reductions \& Make-Whole.} Define caps on any authorized reduction to optional coverage or Pool settlement claims and time-boxed make-whole plans from future fees or rakes, with transparent accounting. Do not describe these reductions as changes to a Holder's underlying Voucher rights.

\textbf{Recoveries.} Any clawback or recovery must be authorized by applicable terms and law, supported by the required process and evidence, and subject to mandatory rights.

\textbf{Reporting.} Publish public post-mortem, remediation timeline, and parameter changes (limits, fees, routes).

\textbf{Important:} Insurance coverage is \textbf{limited}. Some incidents receive no payout after caps are reached; see the Loss Waterfall and exclusions below.

\textbf{When Shared Insurance Will \emph{Not} Make You Whole.}

The Insurance Fund \textbf{does not} cover: (a) redemptions outside the published SLA or venues; (b) losses or coverage reductions beyond policy caps; (c) losses from using delisted or denied routes; (d) fraudulent claims or missing evidence; or (e) jurisdictions where payout is restricted. Payouts, if any, follow the applicable coverage terms and may be \textbf{zero} after caps are reached.

\subsection*{Illustrative Make-Whole Schedule (only if adopted and published):}
\addcontentsline{toc}{subsection}{Illustrative Make-Whole Schedule (only if adopted and published):}

\begin{enumerate}[leftmargin=*]
\item Claims are paid in this order: (a) issuer/guarantor bonds $\to$ (b) pool reserves $\to$ (c) network insurance.
\item If a covered shortfall remains, apply only the reduction to the optional insurance payout or Pool settlement claim that the pre-existing coverage terms and applicable law authorize, up to H\_cap per incident (see Appendix D).
\item Haircut recovery plan:
\end{enumerate}
\begin{itemize}[leftmargin=*]
\item - 25\% of recovered value applied monthly to the covered shortfall until made whole OR
\item - 12-month maximum make-whole horizon; any remaining covered shortfall becomes a recorded loss with a public post-mortem.
\end{itemize}
\begin{enumerate}[leftmargin=*]
\item Every claim produces a receipt: incident ID, affected claim and related Vouchers, coverage-reduction percentage, recovery plan, and appeal window.
\end{enumerate}
\subsection*{11.3 Guarantor Framework}
\addcontentsline{toc}{subsection}{11.3 Guarantor Framework}

This section clarifies who may guarantee what (issuer, Pool, or third-party guarantor), defines the collateral or bonding instruments behind those guarantees, and proposes standard triggers and payout paths. Pools can compete on trust, terms, and expressly offered guarantees without implying that the network automatically guarantees every Voucher.

\subsection*{Baseline Issuer Responsibility (Gift-Card Analogy).}
\addcontentsline{toc}{subsection}{Baseline Issuer Responsibility (Gift-Card Analogy).}

\begin{itemize}[leftmargin=*]
\item Each Voucher is first and foremost the issuer's responsibility: the issuer commits to deliver the specified good or service, or declared cash-equivalent where lawful, under the published redemption terms.
\item Issuers must publish clear terms (who/what/where/when/proof) and disclose limits, venues, and dispute hooks in voucher metadata.
\item If an issuer fails to fulfill, they are the primary party in default; pool or network protections (if any) are secondary layers.
\end{itemize}
\subsection*{Pool-Level Guarantees (Optional, Competitive, Disclosed).}
\addcontentsline{toc}{subsection}{Pool-Level Guarantees (Optional, Competitive, Disclosed).}

A pool may choose to add extra guarantees to vouchers it lists. These are not automatic; they must be explicitly declared in pool metadata and surfaced in receipts.

\subsection*{Common guarantee types:}
\addcontentsline{toc}{subsection}{Common guarantee types:}

\begin{enumerate}[leftmargin=*]
\item Cash-Back Guarantee (Make-Whole in Stable/Reserve Asset)
\begin{itemize}[leftmargin=*]
\item If the issuer defaults or breaches SLA, the pool pays out a defined amount in a designated reserve asset (e.g., stablecoin) up to a policy-capped limit.
\item Funding source: pool-level reserve buffers and/or posted guarantor bonds.
\end{itemize}
\item Swap-Back Guarantee (Reversal / Exit Window)
\begin{itemize}[leftmargin=*]
\item If a voucher cannot be redeemed under declared terms, the pool offers a time-boxed swap-back path (e.g., swap back into the prior asset, or into an approved reserve asset), subject to caps and inventory.
\item This is a liquidity protection, not a promise that every swap is always reversible: it is limited by published caps, windows, and reserve ratios.
\end{itemize}
\item Alternative-Fulfillment Guarantee (Multi-Venue / Substitute Delivery)
\begin{itemize}[leftmargin=*]
\item The pool guarantees fulfillment by routing redemption to an alternative approved provider (e.g., another vetted taxi operator) when the original issuer fails, within a capped quantity/value.
\item This is especially useful for essential services (food/transport) where continuity matters.
\end{itemize}
\item Price/Index Band Guarantee (Optional)
\begin{itemize}[leftmargin=*]
\item For selected Voucher classes, a Pool may commit to keep Swap-out value within a published band relative to its Value Index. If the band is broken, only the coverage adjustment or Swap-back remedy authorized by the pre-existing guarantee terms and applicable law would apply; the guarantee does not itself reduce the Issuer's underlying Voucher obligation.
\end{itemize}
\end{enumerate}
\subsection*{Guarantors \& Bonds (Who can guarantee).}
\addcontentsline{toc}{subsection}{Guarantors \& Bonds (Who can guarantee).}

\begin{itemize}[leftmargin=*]
\item Issuer Bond: collateral posted by the issuer (or locked reserve) that can be drawn down upon verified default.
\item Pool Reserve: pool-owned buffers funded by a portion of pool fees, used for payouts under the pool's advertised guarantees.
\item Third-Party Guarantor Bond: collateral posted by external guarantors (individuals, institutions, insurers, community orgs) that back specific issuers, voucher classes, or the pool as a whole.
\item Note: Guarantor participation is governed by published eligibility criteria, bond sizing, concentration limits, and slashing rules.
\end{itemize}
\subsection*{Triggers (When a guarantee can be claimed).}
\addcontentsline{toc}{subsection}{Triggers (When a guarantee can be claimed).}

Claims must be based on explicit, auditable triggers, such as:

\begin{itemize}[leftmargin=*]
\item Redemption SLA breach (target/max exceeded) with proof of attempted redemption;
\item Verified issuer non-fulfillment or insolvency (as defined by pool policy);
\item Bridge/escrow failure impacting redeemability (when applicable);
\item Governance-declared incident state (emergency pause / run conditions).
\end{itemize}
\textbf{Claim Process (Human-readable and auditable)}.

\begin{itemize}[leftmargin=*]
\item Ticket: user opens a redemption/claim ticket referencing the voucher + proof (QR receipt, ticket \#, required ID type).
\item Verification: pool (or delegated claims module) checks voucher terms, redemption attempt evidence, and issuer response window.
\item Decision: approve/deny within a published dispute window; all outcomes logged.
\item Payout: execute per the published payout path (below), with receipts referencing guarantee type and cap.
\end{itemize}
\subsection*{Payout Path \& Recovery (Aligned to the loss (insurance) waterfall).}
\addcontentsline{toc}{subsection}{Payout Path \& Recovery (Aligned to the loss (insurance) waterfall).}

Guarantee payouts follow a transparent waterfall:

(1) Responsible issuer bond or guarantor stakes $\to$ (2) Pool-level reserves $\to$ (3) Network Insurance Fund, if covered by published policy $\to$ (4) a policy-capped temporary reduction to an optional insurance payout or Pool settlement claim, only if authorized as described in Section 11.2 $\to$ (5) lawful recoveries for proven fraud or abuse.

Recovery proceeds (from issuer settlement, arbitration awards, or legal enforcement) refill bonds/reserves per policy before CLC Pool swap access.

\subsection*{Parameterization (What must be declared).}
\addcontentsline{toc}{subsection}{Parameterization (What must be declared).}

For every pool and voucher class, publish:

\begin{itemize}[leftmargin=*]
\item Guarantor Requirement: None / Recommended / Required.
\item Bond sizing: minimum bond, scaling rule (e.g., \% of issuance or exposure), concentration caps.
\item Guarantee catalog: which guarantee types apply, caps, windows, eligible assets (cash-back asset; swap-back asset).
\item SLA: target/max and claim windows.
\item Disclosures: plain-language ``who is guaranteeing what'', and what is explicitly not guaranteed.
\end{itemize}
\subsection*{Curation Market Principle.}
\addcontentsline{toc}{subsection}{Curation Market Principle.}

Pools are responsible for the guarantees they advertise. A CLC deployment may provide standards, registries, or optional shared insurance policies, but neither CLC nor GEF automatically guarantees Vouchers or Pools. Any coverage must identify the responsible party and be stated in the applicable published policy and Pool terms.

\subsection*{11.4 Anti-Capture Guardrails}
\addcontentsline{toc}{subsection}{11.4 Anti-Capture Guardrails}

Under this template, the following actions would be classified as \textbf{Critical} and would require the adopted highest quorum or approval tier plus a long timelock:

\begin{enumerate}[leftmargin=*]
\item Waterfall structure changes (adding/removing destinations; changing insurance/core ops priority),
\item Registry root changes (canonical voucher/pool registries),
\item Insurance policy scope, coverage-reduction caps, and claims authority changes,
\item Emergency pause scope expansions,
\item Any change that weakens forkability, transparency, or pool sovereignty guarantees stated in this paper.
\end{enumerate}
\subsection*{11.5 Fork \& Exit Procedure (Credible Exit for Communities and Operators)}
\addcontentsline{toc}{subsection}{11.5 Fork \& Exit Procedure (Credible Exit for Communities and Operators)}

If governance were captured or values drifted materially, communities, Pool Stewards, and operators could seek to exit by forking the network governance layer. Preservation of underlying Pools and Vouchers would depend on the deployed contracts, keys, interfaces, infrastructure, and applicable obligations.

\subsection*{Procedure:}
\addcontentsline{toc}{subsection}{Procedure:}

\begin{enumerate}[leftmargin=*]
\item \textbf{Snapshot:} Export canonical registries (pools, vouchers, indices, limits, fee policies) and publish a signed snapshot hash.
\item \textbf{Redeploy:} Deploy new registry roots, router endpoints, and (if needed) a new Waterfall + Insurance policy contract set under a new governance arrangement.
\item \textbf{Re-register:} Pool stewards opt-in by registering their pool addresses under the new registry root (no need to migrate user-held vouchers).
\item \textbf{Client Re-point:} SDKs/UIs add the new registry root as a selectable network profile; default routing can shift via published governance decisions.
\item \textbf{Bridge Period:} Maintain routing bridges where safe to reduce fragmentation; deny-list toxic routes.
\end{enumerate}
\textbf{Design objective:} Exiting canonical registries should not disable otherwise functional local Pools. Actual continuity depends on the relevant contracts, keys, infrastructure, interfaces, and third-party services; federation remains an opt-in discovery layer.



\section*{12. LP Term Sheet (Non-Binding Outline)}
\addcontentsline{toc}{section}{12. LP Term Sheet (Non-Binding Outline)}

This is an illustrative checklist for a future, separately documented program. It is not an offer, a current App feature, or a promise of liquidity, access, insurance, return, value, or exit.

\begin{itemize}[leftmargin=*]
\item \textbf{Eligible Contributions}: Stablecoins, reserve tokens, approved community vouchers.
\item \textbf{Placement}: Into designated CPP pools or the CLC Pool per published mandate.
\item \textbf{Lock-up / Exit}: Program-specific (e.g., rolling 30--90 days; gates under stress).
\item \textbf{Policy-gated fee-credit (ex-post)}, under published caps/windows.
\item \textbf{Risk Sharing}: Any reduction to optional insurance payouts or Pool settlement claims must be expressly authorized by pre-existing terms and applicable law; disclose loss waterfalls and lawful recovery for proven fraud or abuse. Such a reduction does not itself change an Issuer's Voucher commitment or an on-chain balance.
\item \textbf{Reporting}: Monthly dashboards; quarterly attestation of pool health and insurance reserves.
\item \textbf{Covenants}: Use-of-proceeds constraints; oracle hygiene; limit tiers.
\item \textbf{CLC Eligibility (Impact Seeding):} LPs who seed designated pools may become eligible for vested CLC grants based on measured increases in successful network settlement (see \S{}7.2.3), subject to anti-gaming rules and policy caps.
\end{itemize}


\section*{13. Jargon $\to$ Plain Language (Glossary)}
\addcontentsline{toc}{section}{13. Jargon to Plain Language (Glossary)}

\begin{itemize}[leftmargin=*]
\item \textbf{CPP}: A set of contracts that list vouchers, set values, limit swaps, charge fees, and safely hold assets.
\item \textbf{Voucher}: A digital claim for a specific good/service/cash-equivalent.
\item \textbf{Seed (Deposit)}: Add vouchers/tokens into a pool.
\item \textbf{Swap}: Exchange one voucher/asset for another if values and limits allow and inventory exists.
\item \textbf{Value Index}: The pool's pricing table for vouchers vs. a reference unit (may use oracles or governance updates).
\item \textbf{Swap Limiter}: Caps on how much can be swapped over time to prevent runs/arbitrage.
\item \textbf{Router}: Software/contract that finds multi-pool paths for a desired exchange.
\item \textbf{Inventory}: What the vault currently holds and can pay out.
\item \textbf{Guarantor}: A party that stakes collateral to back a voucher/pool against default.
\item \textbf{Redemption SLA}: Expected time to receive the good/service/cash when redeeming a voucher.
\item \textbf{Clearing House (CLC Pool)}: A proposed network-level Pool that may collect fees, hold assets, fund LP programs, or support expressly adopted insurance policies. These functions are not automatic.
\item \textbf{Fee-credit (budget-exit):} A time-bounded, policy-capped authorization (typically via sCLC) to swap fee assets out of designated fee-holding vaults after the Waterfall. It may be set to zero and is not a dividend, yield, or profit-share.
\item \textbf{Rebalancing / Netting Run}: A batch process that executes multilateral cycles/chains across opted-in pools to reduce inventory imbalances and increase successful settlement throughput, subject to published caps and policies.
\item \textbf{On/Off-Ramp:} A regulated service that converts fiat $\leftrightarrow$ approved stable cash-equivalents used to seed or exit pools (e.g., bank transfer, e-money/payment institutions, card-based cash-out), subject to jurisdictional compliance.
\end{itemize}


\section*{14. KPIs \& Health Indicators}
\addcontentsline{toc}{section}{14. KPIs \& Health Indicators}

\begin{itemize}[leftmargin=*]
\item \textbf{Fulfillment Rate}\&\textbf{Redemption Latency} (SLA adherence).
\item \textbf{Reserve Adequacy} by voucher and pool.
\item \textbf{Limit Utilization}\& run-incident avoidance.
\item \textbf{Routing Pass/Fail}\& average hop count.
\item \textbf{Guarantor Performance}\& recovery percentages.
\item \textbf{Protocol Revenue}: pool fees, routing, network rake.
\item \textbf{Net Fee Outcome} (annualized \textbf{sCLC} fee-access, ex-post)
\item \textbf{Governance Responsiveness}: time-to-alarm, time-to-pause, timelock adherence.
\item \textbf{Consumer Protection:} Mis-selling complaints per 1,000 redemptions; median time-to-remedy; delisting decisions per quarter (with public reasons).
\item \textbf{Inventory Skew Index:} per voucher class, dispersion of inventory across pools vs. target bands.
\item \textbf{Rebalance Success Rate:} executed rebalance cycles / attempted cycles; median time-to-rebalance.
\item \textbf{Netting Yield:} (gross routed value - net external liquidity injected) / gross routed value, computed over rebalance windows (method published and timelocked).
\item \textbf{Well-being Outcomes} (profolio program-scoped):
\begin{itemize}[leftmargin=*]
\item Basic-needs coverage proxy (food/transport/health voucher redemption success in target communities).
\item Household resilience proxy (repeat redemption without increased delinquency).
\end{itemize}
\item \textbf{Planetary Regeneration}:
\begin{itemize}[leftmargin=*]
\item Verified ecological outcomes per voucher class (method + auditor published).
\item ``Do-no-harm'' exceptions count (attempted listings rejected by prohibited-list policy).
\end{itemize}
\end{itemize}


\section*{15. Roadmap (Indicative)}
\addcontentsline{toc}{section}{15. Roadmap (Indicative)}

\subsection*{15.1 Current Public Service}
\addcontentsline{toc}{subsection}{15.1 Current Public Service}

As of 24 September 2026, GEF operates the progressive web app at \texttt{cosmolocal.credit} on Gnosis Chain. The App provides account and Wallet access, a public market, and interfaces for creating or using Tokens, Vouchers, Commitment Pools, and direct Swaps. Documentation is published at \texttt{docs.cosmolocal.credit}. Specific availability depends on the App, contracts, inventory, configured limits and fees, network state, and the terms published by issuers and Pool Stewards.

\subsection*{15.2 Future and Optional Work}
\addcontentsline{toc}{subsection}{15.2 Future and Optional Work}

The following sequence is indicative. Its labels describe design phases, not the current protocol release number, delivery dates, or commitments that any feature will be launched:

\begin{itemize}[leftmargin=*]
\item \textbf{v1}: CLC token launch; CLC Pool; fee adapters; governance MVP (quorum + timelocks).
\item \textbf{v1.1: Router SDK \& registry APIs; health dashboards; Insurance Fund policy v1; opt-in rebalance intents + batch netting (cycle-finding) MVP.}
\item \textbf{v1.2}: Cross-domain routing via HTLC/escrow; guarantor module; tiered limit presets.
\item \textbf{v2: Retail on/off-ramps (via regulated partners), personal micro-pools;} compliance plugin marketplace; third-party audits of voucher classes.
\begin{itemize}[leftmargin=*]
\item \textbf{- On-ramps:} convert fiat $\to$ approved stable cash-equivalents that can seed designated pools.
\item \textbf{Off-ramps:} convert approved cash-equivalent stables $\to$ fiat via an approved off-ramp list (bank transfer, e-money issuers, and card-issuing/payment processors on major card rails), with jurisdictional KYC/attestation and geofencing where required.
\item \textbf{Principle:} CLC operators and CPP pools do not operate fiat rails; ramps are provided by licensed third parties under local regulation.
\end{itemize}
\end{itemize}
Multi-hop routing, shared insurance, the proposed CLC/stCLC/sCLC governance system, batch netting, and fiat on/off-ramps remain future or deployment-specific unless the App and applicable published terms expressly state otherwise.



\section*{16. Values \& Evaluation Template (for Listings/Liquidity Mandates)}
\addcontentsline{toc}{section}{16. Values \& Evaluation Template (for Listings/Liquidity Mandates)}

\subsection*{Listing / Mandate Evaluation Rubric (required fields):}
\addcontentsline{toc}{subsection}{Listing / Mandate Evaluation Rubric (required fields):}

\begin{enumerate}[leftmargin=*]
\item Observation: objective facts (issuer history, redemption terms, reserve plan, guarantors, limits).
\item Valuation: what matters (care, risk, capacity, non-dominance) + who is affected.
\item Needs/Boundaries: ecological + social + trust constraints (including prohibited list screening).
\item Request: concrete next step with owner + deadline (approve / revise / reject / human review).
\end{enumerate}
Output must include: risk tier, reserve floor, limiter caps, remedy pathway, and review date.

\textbf{Guardrails}: Timelocked index edits; quorumed limit changes; emergency pause criteria; auditable logs; forkability.

\textbf{Health Indicators}: fulfillment rate; reserve alarms; limit utilization; routing pass/fail; guarantor payout performance.



\section*{17. Legal \& Compliance Note}
\addcontentsline{toc}{section}{17. Legal \& Compliance Note}

CPP coordinates Tokens, Vouchers, Pools, and Swaps whose legal treatment depends on their design, marketing, use, responsible parties, and jurisdiction. A Voucher may be treated as a contractual claim, prepaid service, gift card, payment instrument, credit arrangement, security, taxable supply, or another regulated product in a particular context. Nothing in this White Paper is legal, tax, investment, credit, or financial advice. Issuers, Pool Stewards, service providers, and users must determine and comply with the laws that apply to them before acting.

Use of the public PWA is governed by the \href{https://docs.cosmolocal.credit/governance/terms}{Terms of Service}. Grassroots Economics Foundation (GEF) operates that App and supporting infrastructure. Unless GEF expressly assumes a role in transaction-specific terms, it is not an issuer, Pool Steward, custodian, lender, borrower, broker, guarantor, redeemer, adviser, insurer, or party to obligations between users. Neither GEF nor the protocol guarantees legality, value, liquidity, convertibility, redemption, or performance.

Where fiat on/off-ramps are offered, they must be provided under the responsible provider's own terms and applicable law. The presence of a connector in a CLC interface does not mean that GEF or a Pool operates the underlying banking, e-money, card, or money-transmission service.

\subsection*{17.1 Voucher Class Matrix (Policy Template)}
\addcontentsline{toc}{subsection}{17.1 Voucher Class Matrix (Policy Template)}

This illustrative matrix can help a Pool disclose how it lists and monitors Voucher classes. It is not a legal classification, default guarantee, or representation that every field is implemented by the App.

\begin{longtable}{P{0.30\linewidth} P{0.62\linewidth}}
\toprule
Field & Required disclosure or policy choice \\
\midrule
\endhead
Class Name & For example: food, transport, labor, storage, cash-equivalent, community service, or equipment use. \\
Legal Treatment & The classification reached by the responsible issuer or Steward after jurisdiction-specific review; include material restrictions and required registrations or approvals. \\
Redemption Terms & Issuer identity, offering, quantity or supply, valuation basis, capacity limits, expiry, geography, timing, procedure, fees, restrictions, and material-change process. \\
Identity / Attestation & None, attestation, or identity verification, as lawfully required for the relevant class, value, user, or jurisdiction. \\
Fees & All issuer, Pool, protocol, network, and third-party fees displayed before the relevant action; there is no universal fee tier. \\
Certification / Provenance & If used: source, methodology, scope, expiry, revocation, transferability, and effect on listing or risk treatment. \\
Index Source & Static schedule, OracleQuoter, or governance-updated source; publish method, authority, cadence, bounds, and failure mode. \\
Limits & Per-Voucher, per-account, global, or rolling-window limits, if configured; disclose how limits can change or pause activity. \\
Guarantee / Reserve & None unless expressly offered. If offered, identify the responsible party, assets, funding, caps, triggers, exclusions, evidence, claim process, and applicable terms. \\
Fiat / Stablecoin Provider & Identify the independent provider, supported jurisdictions, eligibility, fees, custody model, and compliance requirements. \\
Disclosures & Plain-language issuer, offering, redemption, risk, complaint, and remedy information displayed before acquisition where required. \\
\bottomrule
\end{longtable}

Pools should declare the class used, enforce their published settings, disclose conflicts and material changes, and avoid representing optional controls as universal protections.

\subsection*{17.2 Minimum Voucher Information}
\addcontentsline{toc}{subsection}{17.2 Minimum Voucher Information}

An issuer should publish, and keep current:

\begin{itemize}[leftmargin=*]
\item \textbf{who:} issuer name, contact details, and responsible legal person or organization;
\item \textbf{what:} the good, service, benefit, or other offering; unit, quantity, supply, valuation basis, and capacity;
\item \textbf{where:} redemption locations, service area, and lawful geographic restrictions;
\item \textbf{when:} issue date, expiry, operating times, fulfillment timing, and any claim window;
\item \textbf{how:} redemption steps and acceptable evidence;
\item \textbf{cost:} issuer, Pool, protocol, network, and third-party fees;
\item \textbf{limits:} per-user, per-transaction, supply, or other material limits;
\item \textbf{fallback:} complaint process and any remedy, guarantor, reserve, or payout path that is actually offered; and
\item \textbf{rights:} any permitted copying, adaptation, redistribution, transfer, suspension, or cancellation and the process for material changes.
\end{itemize}
\subsection*{17.3 Vouchers, Pools, Swaps, and Credit}
\addcontentsline{toc}{subsection}{17.3 Vouchers, Pools, Swaps, and Credit}

\textbf{Voucher.} A Voucher records an issuer's published commitment. The issuer---not GEF merely because it operates the App---is responsible for honoring that commitment. A Holder must assess the issuer, redemption terms, restrictions, expiry, and risks. There is no platform-wide cash-out, stable value, liquidity, or redemption guarantee.

\textbf{Pool.} A Commitment Pool is a contract system governed by its Pool Steward's published rules for admissions, listed assets, valuation, fees, limits, inventory, pauses, configuration, contributions, provenance, conflicts, and any reserve or guarantee. The Steward is responsible for those rules and for any protection it expressly advertises. Listing does not mean that GEF endorses, values, insures, or guarantees an asset.

\textbf{Ordinary Swap.} A Swap is governed by the displayed quote and bounds, Pool rules, smart contracts, network conditions, and applicable law. Asset direction alone does not make one party a lender or borrower, and an ordinary Swap or Voucher redemption does not by itself create, transfer, repay, or discharge a loan.

\textbf{Express credit facility.} Credit mechanics apply only when a Pool or other responsible party expressly presents a transaction as a repayable Advance or Swap Loan and supplies transaction-specific terms. Those terms must identify the parties and disclose the advance and repayment amounts, due date, interest or fees, collateral use, changes in the creditor or Holder, assignment, partial settlement, redemption at any stated Full Value, evidence of discharge, default, and lawful remedies. No network-wide interest rate, repayment period, or redemption-to-repayment rule is implied.

\subsection*{17.4 Access, Risk, and Local Law}
\addcontentsline{toc}{subsection}{17.4 Access, Risk, and Local Law}

Features or asset classes may be restricted, paused, or unavailable based on location, sanctions, eligibility, law, contract state, inventory, limits, security concerns, or third-party services. Identity checks, attestations, consumer disclosures, tax treatment, employment rules, financial-services permissions, or other safeguards may be required. Interface suspension or delisting does not necessarily erase or reverse public-blockchain transactions, balances, contracts, or third-party copies.

Users remain responsible for Wallet credentials, transaction review, taxes, and the legality of their own issuance, Pool operation, acquisition, redemption, and other activity. Smart-contract defects, network congestion or reorganization, OracleQuoter failures, malicious or mistaken configuration, key compromise, issuer non-performance, illiquidity, stablecoin or provider failure, and regulatory change can cause delay or permanent loss. Optional limits, reserves, guarantees, insurance, audits, and governance controls may reduce selected risks but do not eliminate them.



\section*{18. Conclusion}
\addcontentsline{toc}{section}{18. Conclusion}

CLC offers a framework for connecting local production, mutual aid, and expressly documented credit through interoperable Commitment Pools with published registries, value indices, Swap limits, fees, and accountable stewardship. The live PWA provides practical tools for Tokens, Vouchers, Pools, and direct Swaps; the broader routing, governance-token, shared-insurance, and liquidity-coordination mechanisms in this paper remain optional or proposed unless expressly deployed and documented. Transparency and configured controls can improve accountability, but they do not guarantee value, liquidity, redemption, legal status, or protection from loss.


\section*{A. Math Box - Settlement, Fees, Fee Pooling}
\addcontentsline{toc}{section}{A. Math Box - Settlement, Fees, Fee Pooling}

Definitions (per voucher j, per period):

D\_j  := outstanding Voucher commitments valued in the network index (not necessarily debt or a loan)

S\_j  := value of redemptions (settlements) routed through pools

V\_j  := settlement velocity for voucher j

Pool- and Network-Level:

\[
D_j = \sum_k D_{j,k}
\]

\[
S_j = \sum_k S_{j,k}
\]

\[
V_j = \frac{S_j}{D_j}\quad \text{if }D_j>0,\qquad V_j=0\quad \text{if }D_j=0
\]

Aggregate across all vouchers:

\[
D_{\mathrm{tot}} = \sum_j D_j
\]

\[
V_{\mathrm{settle}} = \frac{\sum_j S_j}{D_{\mathrm{tot}}} = \frac{\text{settlement flow}}{\text{outstanding stock}}
\]

Fee volume per period:

\[
F \approx \tau \cdot V_{\mathrm{settle}} \cdot D_{\mathrm{tot}}
\]

Cash-usable fee revenue (convertibility constraint). Let $\chi$ $\in$ [0,1] be the fraction of fee inflows that are cash-eligible/convertible (E\_cash) after slippage and policy constraints. Define:

\[
F_{\mathrm{cash}} \approx \chi \cdot F
\]

Operational break-even and any non-zero sCLC fee-access budget must be evaluated on F\_cash.

LP fee ex-post metrics (per period, rough):

\[
\mathrm{ExPostMetric}_{LP} \approx \frac{\phi \cdot F}{K}
\]

where $\tau$ = average network fee rate; $\phi$ = fee share to LPs; K = LP capital staked.


\section*{B. Illustrative Future Fee Waterfall}
\addcontentsline{toc}{section}{B. Illustrative Future Fee Waterfall}

This is an illustrative design for a future deployment. It applies only if implemented in contracts and adopted through published governance policy; it does not describe a guaranteed current fee or insurance arrangement.

Let F\_in be all fees collected across pools/routers during the epoch.

Eligibility \& conversion note. F\_in may include both cash-eligible fee assets (E\_cash) and in-kind fee assets (E\_kind). The protocol may convert allowlisted fungible assets (E\_cash) into stables/fiat when needed to meet insurance payouts, maintain off-ramp liquidity, and fund core operations---while prioritizing in-network settlement and using liquidity mandates/CLC Pool inventories to reduce settlement latency.

\begin{enumerate}[leftmargin=*]
\item Insurance Reserve Target (IRT): top-up InsuranceFund to Target = $\Sigma$\_p (RW\_p $\cdot$ D\_p),
\begin{enumerate}[leftmargin=*]
\item where D\_p is the pool's outstanding obligations stock (valued in the network index),
\item with RW\_p (risk weight) = f(fulfillment\_rate, issuer\_concentration, limit\_utilization, SLA latency).
\item Allocation = min(IRT - InsuranceFund, MaxTopUp).
\end{enumerate}
\item Core Operations: fixed budget B\_core (timelocked; $\pm$20\% with quorum Q2).
\item Liquidity Mandates: allocate L to approved pools/routers per mandate schedule.
\item Pooled Fees --- allocation of remaining F\_in - (1 + 2 + 3):
\begin{itemize}[leftmargin=*]
\item Protocol Operations/Insurance (non-distributive): $\alpha$
\item Liquidity Programs (incentives/rebates): $\beta$
\item Insurance Buffer (overflow reserve): $\gamma$
\end{itemize}
\end{enumerate}
subject to $\alpha$ + $\beta$ + $\gamma$ = 1, policy bounds, and per-budget caps.

\textbf{Guardrail:} Waterfall allocations are for (i) insurance adequacy, (ii) operations, and (iii) liquidity needed for settlement. They must not be framed or executed as price-support operations.

Under this illustrative design, any CLC acquired through an enabled DEX Float Reduction program would be \textbf{retired} or placed in a disclosed non-voting sink; it would not be distributed to stakers.


\section*{C. KPI Definitions}
\addcontentsline{toc}{section}{C. KPI Definitions}

\begin{itemize}[leftmargin=*]
\item Fulfillment Rate (pool p, period t):  FR\_\{p, t\} = Settlements\_\{p, t\} / RedemptionsRequested\_\{p, t\}
\item Redemption Latency (SLA): median time(issue$\to$redeem) with 90th percentile cap
\item Reserve Adequacy (voucher v): RA\_\{v\} = Vault\_\{v\} / RequiredReserves\_\{v\} (policy function)
\item Limit Utilization (voucher v): LU\_\{v, t\} = Usage\_\{v, t\} / Cap\_\{v, t\}
\item Routing Pass Rate: RPR\_t = SuccessfulRoutes\_t / AttemptedRoutes\_t
\item Avg Hop Count: H\_t = ($\Sigma$ routes hop\_count) / routes
\item Guarantor Recovery: GR\_t = RecoveredFromBonds\_t / ClaimsPaid\_t
\item Protocol Revenue: PR\_t = PoolFees\_t + RoutingFees\_t + NetworkRake\_t
\item Net LP Credit Access (program x): CA\_x = (FeesToLP\_x - AuthorizedCoverageReductions\_x) / AvgStake\_x (annualized)
\item Governance Timeliness: TTA = median(time$\to$alarm), TTP = median(time$\to$pause), TLK = timelock adherence \%
\end{itemize}

\section*{D. Launch Parameters}
\addcontentsline{toc}{section}{D. Launch Parameters}

These values are illustrative proposals for a future governance deployment. Actual values, if any, must be separately adopted, disclosed, and enforced by the deployed contracts; they are not current App defaults or guarantees.

\begin{itemize}[leftmargin=*]
\item Quorum Tiers (of staked voting power):
\begin{itemize}[leftmargin=*]
\item Q1 (Routine): $\geq$ 4\% quorum, >50\% approval.
\item Q2 (Sensitive): $\geq$ 10\% quorum, $\geq$60\% approval.
\item Q3 (Critical): $\geq$ 20\% quorum, $\geq$66.7\% approval.
\end{itemize}
\item Timelocks (minimum delay before execution):
\begin{itemize}[leftmargin=*]
\item T1: 48 hours (Q1 actions)
\item T2: 7 days (Q2 actions)
\item T3: 30 days (Q3 actions)
\end{itemize}
\item Epoch Cadence: default 7 days (policy-set). Each epoch includes:
\begin{itemize}[leftmargin=*]
\item (i) gauge voting window,
\item (ii) emissions calculation,
\item (iii) publication of F\_epoch (may be zero),
\item (iv) sCLC swap windows (if enabled).
\end{itemize}
\item Gauges: canonical list of eligible pools/mandates for incentive direction; edits are timelocked and require Q2 quorum or higher.
\item Emissions Budget (sCLC): a bounded, policy-set maximum per epoch; may be zero; cannot override Waterfall priorities.
\item Anti-gaming: wash-loop exclusion, beneficial-owner clustering, per-entity caps, delayed finalization, and dispute/appeal windows for manipulated metrics.
\item Emergency Pause: immediate (multisig/emergency role), auto-sunsets in 72 hours unless ratified by Q2.
\item Network Fee $\tau$: 20--60 bps (default 30 bps) on routed value
\item Pool Fee Range (steward-set): 0\%--20\% depending on voucher class and risk tier (disclosed on-chain per pool).
\item Network Rake Share r: CLC receives r\% of each pool's collected fees (default 20\%; policy-bounded per pool class).
\item Effective network fee rate per pool: $\tau$\_p = f\_p $\cdot$ r\_p (rake-on-rake).
\item Fee Asset Eligibility Sets:
\item E\_cash (cash-eligible): allowlisted stables/cash-equivalents and (optionally) major liquid tokens.
\begin{itemize}[leftmargin=*]
\item E\_kind (in-kind): non-fiat-redeemable vouchers and other non-convertible fee assets.
\end{itemize}
\item Conversion Policy (fungible assets only): venue allowlists, TWAP windows, max slippage, and monthly caps; quarterly reporting of $\chi$ (cash-eligible share).
\item F\_epoch enablement rule: publish F\_epoch only if (i) InsuranceFund $\geq$ Target and (ii) B\_core is fully funded for the epoch; otherwise F\_epoch = 0.
\item Router Fee Cap: $\leq$ 20 bps per route
\item Reserve Floors by Class: Cash-Equivalent Stable $\geq$ 100\% off-ramp attestations
\item Segregated Risk Lanes:
\begin{itemize}[leftmargin=*]
\item Cash-Equivalent lanes do not cross-subsidize goods/services voucher losses.
\item No routing from Cash-Equivalent to higher-risk classes unless explicitly opted-in per account/pool.
\end{itemize}
\item Default router policy: ``safe-by-default'' (deny cross-class risk unless allowlisted).
\item Optional Coverage-Reduction Cap (per incident): an illustrative maximum of 10\% of the covered claim, with a make-whole schedule, only where pre-existing terms and applicable law authorize it. This does not itself reduce the underlying Voucher commitment or on-chain balance.
\item \textbf{DEX Float Reduction Parameters (if enabled):}
\begin{itemize}[leftmargin=*]
\item \textbf{DEX float definition:} sum of CLC balances in \textbf{allowlisted} external liquidity venues (list on-chain), measured by oracle/indexer method M.
\item \textbf{Trigger:} activate only if DEX float > \textbf{X\%} for \textbf{Y} consecutive days.
\item \textbf{Spend/volume caps:}$\leq$\textbf{Z\%} of trailing-month fee inflows \textbf{and/or}$\leq$\textbf{W\%} of trailing-day DEX volume (use the stricter cap).
\item \textbf{Execution controls:} TWAP window \textbf{T}; max slippage \textbf{S}; randomized delay range \textbf{R}; \textbf{no publication of exact timing} beyond the standing policy.
\item \textbf{Emergency stop:} auto-disable if \textbf{InsuranceFund / Target < I\_min} or if any \textbf{Reserve Floor} is breached.
\item \textbf{Disclosure requirement:}``Not intended to support price; settlement does not depend on DEX price.''
\item \textbf{Receipt transparency:} all executions emit on-chain receipts and are summarized in the epoch report.
\end{itemize}
\end{itemize}

\section*{E. Worked Example - Rake-on-Rake, Convertibility, and Break-Even Runway}
\addcontentsline{toc}{section}{E. Worked Example - Rake-on-Rake, Convertibility, and Break-Even Runway}

Fees are collected in the asset that moves through pools. Some fee assets are cash eligible or convertible (E\_cash), others are in-kind (E\_kind). Define $\chi$ as the trailing-month share of fee inflows that are cash eligible or convertible after slippage and policy constraints.

Example A.

Assume:

\begin{itemize}[leftmargin=*]
\item Average pool usage fee f = 2.00\%
\item Network rake share r = 20\% of pool fees
\item Effective network fee rate $\tau$ = f $\cdot$ r = 0.40\% = 40 bps
\item Monthly cash-denominated requirement B\_cash = Core Ops + required Insurance Top-ups (USD)
\item Cash-eligible share $\chi$ (0 to 1)
\end{itemize}
Then required monthly routed value (in USD-indexed terms) to reach cash break-even is approximately:

\[
R_{\mathrm{required}} \approx \frac{B_{\mathrm{cash}}}{\tau \cdot \chi}
\]

Illustration ($\tau$ = 40 bps):

\begin{itemize}[leftmargin=*]
\item If B\_cash = \$150,000/month and $\chi$ = 25\% $\to$ R\_required $\approx$ \$150,000 / (0.004 $\cdot$ 0.25) $\approx$ \$150,000,000 / month
\item If B\_cash = \$150,000/month and $\chi$ = 50\% $\to$ R\_required $\approx$ \$75,000,000 / month
\item If B\_cash = \$150,000/month and $\chi$ = 100\% $\to$ R\_required $\approx$ \$37,500,000 / month
\end{itemize}
Runway to non-zero sCLC budgets.

Because sCLC fee-access budgets (F\_epoch) are downstream of: (1) Insurance targets, (2) Core Ops, and (3) Liquidity Mandates, F\_epoch may remain zero until R\_monthly consistently exceeds R\_required under conservative $\chi$. This is a long time-frame commitment.

Time-to-target worksheet (update quarterly).

Let R\_0 be current monthly routed value and g be monthly growth rate (e.g., 5\% or 10\%).

Time (months) to reach R\_required is approximately:

\[
t \approx \frac{\log(R_{\mathrm{required}}/R_0)}{\log(1+g)}
\]

Stewards/operators publish quarterly updates of: R\_0, $\chi$, B\_cash, $\tau$ (effective), and the implied t under conservative/base scenarios.


\section*{F. Dataroom Checklist}
\addcontentsline{toc}{section}{F. Dataroom Checklist}

\begin{enumerate}[leftmargin=*]
\item \textbf{Historical Sarafu Network evidence:} archive the earlier service's dated swap volumes, redemption rates, incident history, settlement velocity, aging Vouchers, fulfillment measures, and methodology separately from current Cosmo-Local Credit metrics.
\begin{enumerate}[leftmargin=*]
\item \href{https://dune.com/grassrootseconomics/sarafu-network}{Historical Dune Analytics dashboard}
\end{enumerate}
\item \textbf{Proposed fee-enforcement evidence, if implemented:} prove that the FeeHook contract and registry gating are integrated into the settlement execution path rather than enforced only by an interface.
\begin{enumerate}[leftmargin=*]
\item Show that the applicable Router rejects any hop whose receipt does not prove fees were collected under the adopted policy.
\item Published test vectors + failure cases + list of compliant pools (signed by registry root).
\item (segregated cash-equivalent lanes)
\item \href{https://software.grassecon.org}{Software Description}
\end{enumerate}
\item \textbf{Guarantor Framework pack:} eligibility, collateral, premiums, triggers, dispute/appeal examples; 2--3 sample Vouchers distinguishing issuer responsibility from any expressly offered Pool guarantee or network insurance.
\begin{enumerate}[leftmargin=*]
\item \href{https://sarafu.network/reports}{Historical Sarafu Network reports archive}
\item \href{https://youtu.be/5yGaP31bvs0?si=fZ-AMTEXsmVR8Ulv}{Historical year-in-review presentation for operations in Kenya}
\end{enumerate}
\item \textbf{Legal memo:} jurisdiction strategy, Voucher classes, cash-equivalent policies, KYC/attestation options, geofencing, mis-selling controls, and the ordinary-Swap versus express-credit distinction.
\begin{enumerate}[leftmargin=*]
\item \href{https://docs.cosmolocal.credit/governance/terms}{Current Cosmo-Local Credit Terms of Service}
\end{enumerate}
\item \textbf{Governance hardening:} COI template, recusal mechanism, sanctions ladder, delisting due-process templates, and examples of time-locked index/limit edits.
\item \textbf{Source and assurance pack:} source and licence inventory, available audit reports, invariants, build provenance, and monitoring dashboards.
\begin{enumerate}[leftmargin=*]
\item \href{https://github.com/grassrootseconomics}{GitHub}
\end{enumerate}
\end{enumerate}

\end{document}
